\documentclass{jsarticle}
\usepackage[dvipdfmx]{graphicx}
\usepackage{chemfig}
\usepackage{amsmath}
\usepackage{amssymb}
\begin{document}
\title{大域的微分多様体を \\
単体的積分と単体的微分へと \\
書き換えることにしたレポート}
\author{Masaaki Yamaguchi}
\date{}
\maketitle

   ベータ関数をガンマ関数へと渡すと、
   $$ \beta(p,q) = {{\Gamma(p)\Gamma(q)} \over {\Gamma(p + q)}} $$
   ここで、単体的微分と単体的積分を定義すると
   $$ t = \Gamma(x) $$
   $$ T = \int \Gamma(x) dx $$
   これを使うのに、ベータ関数を単体的微分へと渡すために、
   $$ \beta(p,q) = - \int {1 \over {t^2}}dt $$
   大域的微分変数へとするが、
   $$ T_m = \int \Gamma(x) dx_m $$
   この単体的積分を常微分へとやり直すと、
   $$ T^{'} = {t^{'} \over {{\log t}}}dt + C(\text{Cは積分定数}) $$
   これが、大域的微分多様体の微分変数と証明するためには、
   $$ dx_m = {1 \over {\log x}}dx, dx_m = {(\log x^{-1})}^{'}  $$
   これより、単体的微分が大域的部分積分へと置換できて、
   $$ T^{'} = \int \Gamma(\gamma)^{'} dx_m $$
   微分幾何へと大域的積分と大域的微分多様体が、加群分解の共変微分へと
   書き換えられる。
   $$ \bigoplus T^{\nabla}  = \int T dx_m, \delta(t) = t dx_m  $$


   これらをまとめると、ベータ関数を単体的関数として、
   これを常微分へと解を導くと、単体的微分と単体的積分
   へと定義できて、これより、大域的微分多様体と大域的積分多様体
   へとこの関数が存在していることを証明できるのを上の式たちから
   わかる。この単体的微分と単体的積分から微分幾何が書けることがわかる。

   
   これらは自分が導いた方程式たちとおもうようにする。

   
   微分の本質は、差と商の計算であり、積分の本質は加群である。
   この微分と積分を合わせた計算が加群分解であり、ガロワ群である。


   アカシックレコードにいると思っていた私の子どもとも勘違いしていた、
   子どもはまぼろしであり、架空のストーリーの作為体験の子ともと思う次第
   であり、架空のストーリーに載せられているとの妄想癖での、
   このストーリーに載せられている自体の妄想が作為体験である。
   作為があったとは違う。日常の作為と妄想の作為は違う。

   レドックス電池を言っていたから、荒木先生の経営している病院の
   医者である江松先生が私を助けようとして、
   作為体験のために、私のあのときまでの情報、
   私の周りの情報を使って、架空のストーリーを立てて、作為体験として、
   あのときは、ＭＲＩ室の壁を取っていて、壁が無かった。
   先生は、東北の地震があるのを事前に知っていたとしか思えないことを、
   福山市の私の散歩コースに現れて、挨拶とを交わしていて、
   本当に庇ってくれて申し訳ないです。ＳＲＳ速読法をしていたのを、
   占さんがアルバイトのあとに、トポロジーのレポートを見せていたら、
   気づいたらしく、東北のことを知っていたらしく、史記と三国志を
   私に渡した上で、読書力を中心として、勉強をしなさいとフォローを
   してくれて、していたら、２月２３日に持田さんと伊藤一朗さんが
   ラジオ放送に出ていて、シングルをリリースしましたと言うので、
   YouTubeを見ると、一箇所だけ動画がとんでいて、かばうのを
   彩さんにすればいいのに、統合失調症の真逆のうれしいのに悲しむとは
   感情と現状での置かれている立場から違うが、似たような作用の症状でもあり、
   「存在と無」を啓文社に予約したのが、取りに行くのに１８：００
   にしたのが、私の顛末であり、全部私の責任です。７ヶ月間びっしりと
   いたぶられました。増援は蔵王病院らしく、とっくに気づいています。
   史記は、小方君と占さんのおかげで助かって、三国志が私に増援がくると
   助言していた。占さんが蔵王病院に行かしたのは、１０月１日に彩さん
   がめざましテレビを辞めたからと、お父さんの体調を気にしてのことと
   ３年前くらいからわかっていました。２０１３年に彩さんがめざましテレビ
   を辞めたのを知ったが、２０１８年の６月に何時辞めたかを知り、
   携帯サイトでのKindleから１８７７回目のめざましテレビの放送で
   辞めたのを知り、姉が怒っていたのを、いろんないきさつがあったのを
   知り、私の現状を自身にふりかかっている出来事に怒ることが、
   少なくなりました。グロス・グラスマンは、間所くん、青木くん、
   伊澤くん、圭一さん、寺林さん、鷺坂くん、下神くん、梶岡くん、
   片寄くん、鶴元くん、圭一くんです。
   鶴元くんの元は、元寇の元であり、真逆で、中国の反対の日本でもある。
   占さんを入れたいが、周りの批判を受けるので、入れることができない。
   小方くんは、ニールス・ボーアである。
   大域的微分多様体の基軸になる、基本群をもとめることができたのは、
   虚数が不思議、平方根が不思議と、$$ \pi(\chi,x) = [i\pi(\chi,x), f(x)] $$
   の方程式ができたことの根本の源は、小方くんである。
   一般相対性理論の式は、
   $$ \kappa T^{\mu\nu} = R_{\mu\nu} + {1 \over 2}g_{ij}\Lambda $$
   これを大域的微分多様体と同型といえるのは、
   $$  {d \over df}F = {d \over df}\int\int{1 \over {(x \log x)^2}}dx_m
   + {d \over df}\int\int{1 \over {(y \log y)^{1 \over 2}}}dy_m $$
   ここで、一般相対性理論の多様体積分は、
   $$ \int\left( {R_{\mu\nu} + {{1 \over 2}g_{ij}\Lambda}}\right)\mathrm{
   dvol} $$
   $$ \int {\kappa T^{\mu\nu}}\mathrm{dvol} = - 2 R_{ij} $$
   これは、大域的微分方程式の重力と反重力方程式と同型であり、
   $$ = \int\left(R_{\mu\nu}\right)\mathrm{dvol} + \int{{1 \over 2}g_{ij}
   \Lambda}\mathrm{dvol} = \text{重力場方程式 + 反重力場方程式(斥力)} $$
   $$ = {d \over df}\int\int{1 \over {(x \log x)^2}}dx_m
   + {d \over df}\int\int{1 \over {(y \log y)^{1 \over 2}}}dy_m $$
   と、多様体積分が大域的微分多様体となる。


   そうゆう経緯であり、なぜアカシックレコードの私の子どもが
   ２０２３年の１０月に狙いをさだめているかがわかるのは、
   私がご先祖様にアクセスするのと原理が同じだが、ドアーを使う
   アカシックレコードへのアクセスによっても、微分幾何の量子化から
   私の大域的微分多様体の計算方法が間接的にわかったことでもあり、
   直接でもある、思念同士の交流でもある方法だからです。
   彩さんが今産みたいとおもっても、あの子は待つと思うという次第です。
   １２月２５日に狙いを定めていて、うちのユウと同じことをしようとしている
   という次第です。彩さんは、マスターベーションで産みたいらしいので、
   私は文句を言わんが、あまり、私をいじめないでほしいです。
   あの子は了解済みらしいです。

\title{放射性物質吸収体とRNA干渉としてのオイラーの公式}
\author{Masaaki Yamaguchi}

  　オイラーの公式から、XY平面での曲平面をZ軸での虚軸として、
  ゼータ関数をZ軸にAssembile-D-braneのシリンダーを束ねると、
  DNAのオミクロン曲線へと２重螺旋構造を再現できて、
  素数の分布が双対としてのDNAの父親と母親から受け継いでいる、
  染色体として、この説明でのアミノ基の構造が、ヒッグス場とオイラーの定数
  との加群で作られるガンマ関数における大域的微分多様体として、
  表現できることを、
  $$ \int \Gamma(\gamma)^{'}dx_m = {d \over df}F + \int C dx_m $$
  $$ = e^f + e^{-f} \ge e^{f} - e^{-f} $$
  $$ e^{-\theta} = \cos(i x \log x) - i \sin(i x \log x) $$
  $$ e^{i \theta} = \cos \theta + i \sin \theta $$
  これらの式たちから、DNAの生体地図としての、染色体の構造式は、
  オイラーの公式と同型としての、
  $$ \int \Gamma(\gamma)^{'}dx_m  = 2(\cos (i x \log x) - i \sin (i x \log x))$$
  と表される。

  この式は、反重力が放射性物質から励起される熱エネルギーを
  、カーボンナノチューブとコバルト６０との
  複合体でもある、二量体としての結晶石が、吸収することができる、
  生体エネルギーの人体の保温と同じ原理で、反重力の放射性物質を吸収
  することがこの結晶石が、オイラーの定数の大域的微分多様体から
  導く金属の結晶構造の配置と同じ原理で表される式として、
  ヒッグス場とオイラーの定数の大域的微分多様体の式が、
  この二量体の結晶石の構造式として、求まる。

  これらの式たちは、次の理論で証明できる。
  数値は結果であり、アイデアから式、そして値をもとめることが、
  エレガントな解き方と言える。
  

    ベータ関数をガンマ関数へと渡すと、
   $$ \beta(p,q) = {{\Gamma(p)\Gamma(q)} \over {\Gamma(p + q)}} $$
   ここで、単体的微分と単体的積分を定義すると
   $$ t = \Gamma(x) $$
   $$ T = \int \Gamma(x) dx $$
   これを使うのに、ベータ関数を単体的微分へと渡すために、
   $$ \beta(p,q) = - \int {1 \over {t^2}}dt $$
   大域的微分変数へとするが、
   $$ T_m = \int \Gamma(x) dx_m $$
   この単体的積分を常微分へとやり直すと、
   $$ T^{'} = {t^{'} \over {{\log t}}}dt + C(\text{Cは積分定数}) $$
   これが、大域的微分多様体の微分変数と証明するためには、
   $$ dx_m = {1 \over {\log x}}dx, dx_m = {(\log x^{-1})}^{'}  $$
   これより、単体的微分が大域的部分積分へと置換できて、
   $$ T^{'} = \int \Gamma(\gamma)^{'} dx_m $$
   微分幾何へと大域的積分と大域的微分多様体が、加群分解の共変微分へと
   買い換えられる。
   $$ \bigoplus T^{\nabla}  = \int T dx_m, \delta(t) = t dx_m  $$


   これらをまとめると、ベータ関数を単体的関数として、
   これを常微分へと解を導くと、単体的微分と単体的積分
   へと定義できて、これより、大域的微分多様体と大域的積分多様体
   へとこの関数が存在していることを証明できるのを上の式たちから
   わかる。この単体的微分と単体的積分から微分幾何が書けることがわかる。

   この説明から単体的微分と単体的積分が結晶石の二量体としてと、
   RNA干渉としての、オイラーの公式をも、構造式として表せることを
   理論建てしている。

   アメリカの国防総省のUFOの解体から、リバースエンジニアリングした
   UFOもどきが、この航空機が飛行するのに、機体の中心部に
   放射性物質による結晶石をつかっているのが、宇宙人の発想の逆を言っている
   とわかり、早坂先生の反重力発生装置が解決しているのを、
   この反重力発生装置の熱エネルギーを発する放射性物質自体の
   航空機への放射性物質の被爆を防ぐのに、UFOは中心部へこの放射性物質
   を吸収する結晶石を置いているのと、その放射性物質の
   副次的エネルギーを、UFOの副次的動力源として、
   エネルギーのリサイクルをも、人間が核エネルギーの
   原子力発電所のリサイクルエネルギーと同じ発想で、
   UFOの動力源の補助をもこの放射性物質の結晶石がしているというのが、
   わかった次第でもあります。
   なぜ、このUFOもどきをつくっても少ししか上空しないかは、副次的エネルギー
   としての、主エネルギーに比べて、少ないという、この機知
   の仕組みで歴然としています。

   UFOが地球に頻繁に現れ始めたのが、原爆のあとと、ATR研究所が
   ブレインインターフェースを作り始めたら、宇宙人による人へのインプラント
   が頻繁に始まり、本当に理論と工学のからくりがわかりやすく、
   一人さんの持論を宇宙人までしている始末でもあります。

\begin{document}
\title{Abandust judgement on economy strategiest \\
built with acceptant of movement \\
Money demand on law of judgement}
\author{Masaaki Yamaguchi, and investigate from AYA TAKASHIMA \\
also future from my son being have with \\
quantum level on differential geometries}

高島彩さんの研究論文をアカシックレコードと瞑想法で見ての書き出しとしての
論文

   Retry for future to one result with adapt from adulsent for
aquirance of trainning. For must friends of money are treat on
law of judgement, economy need to treat for aquirered of demands,
this demands supply from rest of  one's suppliement.
This balance end without a closed accidennt of out come for
restorance with supplied of demands exceed.
These system are constructed with Actor, Demand, Supply of
being retrayed from Jones manifold around economy of zone
with accesority, verisity of curve in up, down, eternal
rout of stages.

    This system also monument with partial and economy of
similist from law of judgement in universe and money, economy,
human society, estrand of moment in last eternal stage.
  

   Jones manifold equation built with Euler equation and Euler product,
this concernd with concept from Higgs and partial fields,
moreover this concernd of circle and Euler product from
Global integrate and Volume manifold in Sheap manifold of integrate
with group and toplogy different equation.

$$ e^{- \theta} + e^{i \theta} = \sin (i x \log x) = e^{-f} + e^{f} $$
$$ + \cos (i x \log x) = e^{f} - e^{-f} $$

  自動車における道路での速度とそれに対しての加速度を
  平地と上り坂、下り坂における曲がり角での曲率$R_{ij}$が、
  広中平祐先生の庭園理論と同じテーゼで、彩さんが
  経済理論の景気回復と景気刺激剤、景気恐慌が
  どのようにして、自動車運転と自動車の交通規則と
  同じ理論で作用するかを述べているのが、この論文で説明されている。 


  $$ {d \over dx}\left({d \over d r}R_{ij} = \sin
	  \begin{vmatrix} {l_2 \over {2 \pi}} 
- {l_1 \over {2 \pi}}\end{vmatrix}  < 1 \right) = 0 $$


      Curveの度合い $R_{ij}$、潮汐力の差、${d \over d f}F$とは、
   直接的には違うが、間接的には同じであり、rとRの大域的変数
   として、半径rの多様体が曲率的には$R_{ij}$に大域的微分で作用をしている。
   遠隔的に作用している。


      グリーシャ先生からリッチ・フロー方程式が、大域的微分になるのを
   教わった。広中平祐先生の4重帰納法で、ガンマ関数における大域的
   微分多様体が、オイラーの定数の多様体積分に使えるのがわかった。
   これがヒルベルト空間としてのサーストン多様体が、
   代数多様体としての微分幾何として、ゼータ関数が素数分布として
   展開されるのがわかった。

  $$ R_{ij} = |R_2 - R_1| < 1, |R_2 - R_1| = 0 $$
  $$ R_{ij} = {R_2 \over R_1} < 1, {R_2 \over R_1} = 1 $$
  $$ {d \over dt}g_{ij} = - 2 R_{ij}, {d \over d f} F(x) = 
  \int \Gamma(\gamma)^{'}dx_m = \int \Gamma dx_m + {d \over d \gamma}
  \Gamma \le e^{f} - e^{-f} \le e^{-f} + e^{f} $$

\includegraphics[width=6cm,clip]{curve11}
\includegraphics[width=6cm,clip]{curve112}
\includegraphics[width=6cm,clip]{curve12}
\includegraphics[width=6cm,clip]{curve121}
\includegraphics[width=6cm,clip]{curve1211}
\includegraphics[width=6cm,clip]{curve122}

$ R_2 < R_1 \text{の場合は、}$
左からR1に接近して、R1へ行くが、${R_2 \over R_1} < 1 $
より、曲率$R_{ij} < 1$であるから、平坦の道路では、気をつけるべし、
曲率$R_{ij} = 1$では、CURVEには気をつける。
曲率$R_{ij} > 1$では、$R_{i > j}$であり、iの方で気をつける。
これを経済の日経平均相場で、景気の上がりでは、curveがどう分類されるかで、
気をつけるべき時期がわかり、景気の下がりでは、同じくcurveでどう分類するかを
きめる。平坦のときが一番気をつける、油断できない景気の税と立法での
判断が、自動車の走行分岐での心理作用が、景気判断のどの道路状況で
その場の気をつける心理と同じ作用をする。
経済の需要が供給を下回っている状況では、$\text{demand} \le \text{supply}$
では、需要のために、物価を下げるが、この状況では、$ R_1 < R_2, $
であるので、株は$R_2$の方へ曲率がまがり、このときに、需要の
企業の生産が減ってきて、このときに、金融機関が気をつけないと、
企業が破綻する可能性がある。このシチュエーションが世界恐慌であった。
このように、曲率の場合分けをしていないと、間違った判断で、
曲率に気をつけないから、交通事故が起こるということである。
広中平祐先生の庭園理論と同じ考えのモチーフを高島彩さんは、
研究論文で提出している。

   高島彩さんの成蹊大学での研究論文を参照させてもらえて、
書き出した経済の需要と供給に対しての物価上昇と経済沈滞
においてのインフレとデフレにメスを入れる税と景気対策
のタイミング、このタイミングがJones多項式の周期の終わり
にバランスを崩すと、景気の極率としてのCURVEの潮汐力での
差分率の曲率が１以下になるときの、心理作用を考えるべき
という、経済消費における物価調整の作用をかけるメカニズム
についての論文 　

参考文献: 父、母、彩さん、ナッシュ先生、益川先生、ワインバーグ先生、まどかさん


\title{慣性の法則と、回転体による反重力 \\
タイミングとともに、別ベクトルとして重力を加える法則}
\author{Masaaki Yamaguchi}

   慣性の法則の慣性力が働くときと同時に、
   前もって、タイミングを計り、慣性力とともに、
   逆向きのベクトルとして、慣性力の力と同じエネルギーの
   重力を生成して、慣性力の影響を無くす。この力を反重力という。
   慣性力と逆向きのベクトルとして、反重力として、
   慣性力に加えて、別のベクトルとする。

   反重力の本質は、重力のベクトルの違いであり、
   核エネルギーが特殊相対性理論と一般相対性理論を使って、
   放射性物質吸収体に、この原理を使い、核エネルギーをリサイクルしている。
   これを宇宙人は知っていて、UFOにこの原理を使っている。

   Higgs場は、重力であり、周りを反重力が覆っている。
   一般相対性理論のベクトルの向きを変えた力の装置が、
   ラムダ・ドライバであり、反重力である。

   重力は熱エネルギーである。ベクトルの向きが違う。
   反重力は吸収体である。
   慣性の法則が働いているときに、これが手に入ったが、離れたになり、
   重力をベクトルの向きを変えて、同じエネルギーを加えた場合、
   量子コンピュータが、そのときの差分を使って、反重力が慣性の法則
   を変えて、宙に浮くUFOにしている。 

   まとめると、慣性の法則のときに、重力のベクトルを違う向きとして、
   加えると、手に入ったが、離れたになり、これが反重力である。 

   UFOが現れる前に、ヘリコプターが現れて、そのあとに、UFOが実際に現れる。
   慣性の法則のしばらくして、反重力として、UFOとして、
   宇宙人が人に教えている。

   反重力の生成は、回転体の電磁気力からくるちからであり、ローレンツ力と
   同じでもあり、機体の内部で生成される力であり、そのために、放射性
   物質吸収体を結晶石として、このエネルギーからくる放射力を吸収する
   ために、結晶石に反重力の放射線を吸収する。その上に、外部の
   慣性力を内部の反重力で、手に入ったら、離れたをしている。


\begin{document}
\title{反重力発生装置の各構成物質}
\author{Masaaki Yamaguchi}



  原子振動子　セシウム Cs


  形状記憶合金 (Fe・Co60・Pt・Al)$H_2 SO_4, Al_2 O_3$


  反重力発生器 Pr(パラジウム合金)電磁場生成He,$H_2$Mg,Al(摩擦熱)


  結晶石 Co60,Pt,Prの反陽子


  緩衝剤 慣性力感知器 有機化合物とS化合物(シクロアルカン$C_n H_{2n}$
  方位と位置探知機にも使う


  量子コンピュータの量子素子 Pt,Ag,Au（電子と電流のエネルギー経路)半導体に
  使う
  回路の合成演算子 


  $$ x^y = {1 \over {y^x}} \to $$

  $$ \pi(\chi,x) = [i\pi(\chi,x),f(x)] $$
  $$ {d \over df}F = F^{{f}^{'}} $$


  ディスプレイの電磁迷彩 (Al,Mg) $\to$ S合成子 $Si O_2 $
  プリズム C

\title{Fiber of layzerium in amalugam and stephany of parm \\
in antigravity from electricity \\
and magnity of non condition's power}
\author{Masaaki Yamaguchi}

   大域的多重積分と大域的無限微分多様体を定義すると、
   Global assemble manifold defined with infinity of differential
   fields to determine of definition create from Euler product and
   Gamma function for Beta function being potten from integral manifold
   system. Therefore, this defined from global assemble manifold from
   being Gamma of global equation in topological expalanations.

   $$ \int f(x)dx = \int \Gamma(\gamma)^{'}dx_m $$
   $$ = 2 (\cos (i x \log x) - i \sin(i x \log x)) $$
   $$ \left({{\int f(x) dx} \over \log x}\right) = \lim_{\theta \to \infty}
    \left({{\int f(x) dx} \over \theta}\right) = 0,1 $$
    $$ e^{i \theta} = \cos \theta + i \sin \theta $$
    $$ \left({\int f(x) dx}\right)^{'} = 2 (i \sin(i x \log x) - \cos (i x \log
    x)) $$
    $$ = 2(- \cos(i x \log x) + i \sin(i x \log x)) $$
    $$ (\cos (i x \log x) - i \sin(i x \log x))^{'} $$
    $$ = {d \over {d {e^{i\theta}}}}\left((\cos, - \sin )\cdot
    (\sin, \cos)\right)  $$
    $$ \int \Gamma(\gamma)^{''}dx_m = \left({\int \Gamma(\gamma)^{'}}dx_m
    \right)^{\nabla L} = \left({{\int \Gamma dx_m} \cdot {{d \over d\gamma}
    \Gamma}}\right)^{\nabla L} \le 
    \left({\int \Gamma dx_m + {d \over d\gamma}\Gamma}\right)^{\nabla L} \le 
    \left({e^{f} - e^{-f}} \le {e^{-f} + e^{f}}\right)^{'}  $$
    $$ = 0,1 $$

    Then, the defined of global integral and differential manifold
    from being assembled world lines and partial equation of deprivate
    formula in Homology manifold and gamma function of global topology
    system, this system defined with Shwaltsshild cicle of norm from
    black hole of entropy exsented with result of monotonicity for
    constant of equations.
    以上 より、大域的微分多様体を大域的 2重微分多様体として、
    処理すると、ホモロジー多様体では、種数が１であり、特異点では、種数が０
    と計算されることになる。ガンマ関数の大域的微分多様体では、
    シュバルツシルト半径として計算されうるが、これを大域的2重微分で
    処理すると、ブラックホールの特異点としての解が無になる。

    川島隆太先生が、私みたいな人が、頭の中で偏微分や多重積分、部分積分、
    共変微分や、置換積分や、複素微分や複素積分を考えていると思っていた
    と言っていたが、それ以上の大域的微分や大域的2重微分や、
    大域的多重積分や大域的偏微分や大域的部分積分、微分幾何の量子化、
    単体積分、単体微分などの、大学院生でも考えない、思いもしない、
    実際の式を書くと誰でも、その関数と多様体が存在するなと思うしか
    できない、私の記述式を見ると納得する理論式でもあると、
    私も自慢できる理論であると自負できる。

  $$ \log x|_{g_{ij}}^{\nabla L} = f^{f^{'}}, F^f|_{g_{ij}}^{\nabla L}
  |: x \to y, x^p \to y $$ 
  $$  f(x) = \log x = p \log x, f(y) = p \log x  $$
  大域的偏微分方程式と大域的多重積分は、それぞれ次のように成り立っている。
  It's defined with global partial deprivate formula and assemble manifold. 
  $$ {d^2 \over df^2}F = F^{f^{'}}\cdot f^{f^{''}},  $$
  $$ \int \int F dx_m = F^f \cdot F^{(f)^{'}} $$
  $$ {d \over df dg}(f,g) = (f \cdot g)^{f^{'} + g^{'}} $$
  大域的部分積分も、次のように成り立っている。
  And, global parcial integral equation also established with global
  topology computations. therefore, this circutation exclude for all of
  topology extension ideas.
  $$ (F^f \cdot G^g) = \int (f \cdot g)^{f^{'} + g^{'}} $$
  $$ \int \int F\cdot G dx_m = [F^f \cdot G^g] - 
  \int (f \cdot g)^{f^{'} + g^{'}} $$
  大域的部分積分の計算は、
  $$ \int {d \over {df dg}}FG = \int F^{f^{'}}G + \int F G^{g^{'}} $$
  $$ \int F^{f^{'}}Gdx_m = [F^f G^g] - \int F G^{g^{'}}dx_m$$
  大域的商代数の計算は、
  Global quato algebra computations also,
  $$ \left({F(x) \over G(x)}\right)^{(fg)^{'}} = {F^{f^{'}}G - FG^{g^{'}}
  \over G^{g}} $$
  大域的偏微分方程式は、縮約記号を使うと、

  Global partial manifold explain with reduction of operator for being
  represented from partial equation of references.
  $$ {d \over df dg}FG = {d \over df_m}FG $$
  $$ {\partial \over \partial f_m}FG = F^{f^{\mu\nu}}\cdot G
  + F \cdot G^{g^{\mu\nu}}, \int F dx_m = F^f $$
  多様体による大域的微分と大域的積分が、エントロピー式で統一的に表せられる。

  Entorpy computation from integrate for being manifold of global deprivate
  and integrate formulas.
   ベータ関数をガンマ関数へと渡すと、
   Gamma function sented for Beta function of inspiration with monitonicity
   of component with differential and integral being definitions.
   $$ \beta(p,q) = {{\Gamma(p)\Gamma(q)} \over {\Gamma(p + q)}} $$
   ここで、単体的微分と単体的積分を定義すると
   This area defined with different and integrate of component with
   global topology.
   $$ t = \Gamma(x) $$
   $$ T = \int \Gamma(x) dx $$
   これを使うのに、ベータ関数を単体的微分へと渡すために、
   for this system used for being is conbiniate 
   with beta function for componenet of deprivate eqaution.
   $$ \beta(p,q) = - \int {1 \over {t^2}}dt $$
   大域的微分変数へとするが、
   This point be for global differential variable exchanged,
   $$ T_m = \int \Gamma(x) dx_m $$
   この単体的積分を常微分へとやり直すと、
   This also point be for being retried from ordinary differential computation
   for component of integral manifold.
   $$ T^{'} = {t^{'} \over {{\log t}}}dt + C(\text{Cは積分定数}) $$
   これが、大域的微分多様体の微分変数と証明するためには、
   This exceed of proof being for being defined with deprivate variable of
   global differential formula, this successed for true is,
   $$ dx_m = {1 \over {\log x}}dx, dx_m = {(\log x^{-1})}^{'}  $$
   これより、単体的微分が大域的部分積分へと置換できて、
   Therefore, this exchanged from mononotocity deprivation from global
   parital integral manifold is able to, 
   $$ T^{'} = \int \Gamma(\gamma)^{'} dx_m $$

   微分幾何へと大域的積分と大域的微分多様体が、加群分解の共変微分へと
   書き換えられる。
   After all, these exchanged of monotonicity deprivation successed from
   being catastrophe of summativate of partial and assemble of deprivations
   for differential geometory of quantum level to global integral and
   differential manifold.
   $$ \bigoplus T^{\nabla}  = \int T dx_m, \delta(t) = t dx_m  $$


   this exluded of being conclution are which beta function
   evaluate with mononicity from ordinary differential equation
   be resulted from component of deprivation and integral expalanations.
   This cirtutation be resenbled to define with global topoloty of
   extention extern of deprivate and integral of manifolds
   estourced with quantum level of differential geometry be proofed with
   all of equation anbrabed from Euler product.
   これらをまとめると、ベータ関数を単体的関数として、
   これを常微分へと解を導くと、単体的微分と単体的積分
   へと定義できて、これより、大域的微分多様体と大域的積分多様体
   へとこの関数が存在していることを証明できるのを上の式たちから
   わかる。この単体的微分と単体的積分から微分幾何が書けることがわかる。



   虚数の虚数倍した値が超越数のx倍と同じとすると、
   超越数の$2 \pi m$倍が$n=1$でiとなると定義すると、
   次の式たちが導かれる。　
   Imaginary pole circlate with twigled of pole in step function
   from Naipia number of assembled from equalation of defined are
   escourted to be defined with next equations.


   These defined equation are climbate with idea of equation from
   Caltan of imaginary number of circulation.
   カルタンを超えているアイデアと数式たちでもある。

   $$ i^i = (\sqrt{i})^{\sqrt{i}} = e^{x \log x} $$
   $$ e^x = i, e^{2\pi m} = i^n, {d \over d x}e^{2\pi m} = i^n $$
   $$ e^{i \theta} = \cos \theta + i \sin \theta, e^{2 \pi m} = i^n $$
   $$ 2 \pi m = n, e = i \text{($e=i$となる$e^{i \theta} = \cos \theta 
   + i \sin \theta $
   の範囲で$2\pi m$ がある。)}$$
   $$ {d \over d i}i = i^i = e^{x \log x}, m={1 \over 2 \pi},l=2\pi r, 
   \pi = {l \over 2r} $$
   $$ ||ds^2|| = e^{-2 \pi T||\psi||}[\eta_{\mu\nu} + \bar{h}(x)]
   dx^{\mu}dx^{\nu} + T^2 d^2 \psi $$
   $$ e^x = i, e^{2\pi m} = i^n, e = i $$
   $$ {{[\eta_{\mu\nu} + \bar{h}(x)]dx^{\mu}dx^{\nu}}/i} = H \Psi
   = i \hbar \psi, H \Psi = {1 \over i}[H, \Psi] $$
   ハイゼンベルク方程式が$ AdS_5 $多様体の原子レベルの方程式も表せられる。 
   微分幾何の量子化の式は、
   Hisenbelg equation are represeted from AdS5 manifold to particle level
   equalation of quantum levle of differential geometry entranced.
   $$ \bigoplus {({i \hbar} ^{\nabla})}^{\oplus L} = e^{{i x \log x}^{
	   {e^{i x \log x}}^{'}}} $$
   $$ \Gamma = \int e^{-x}x^{1 - t}dx $$
   $$ \gamma = \int e^{-x}x^{1 - t}\log x dx $$
   $$ = \int \Gamma(\gamma)^{'} dx_m$$
   $$ = {d \over d \gamma}\Gamma $$

   Eight differential geometry are each intersect with own level of
   concept from expalanation of Euler product system.
   This component of three manfold sergeried with geometry of destroy and
   desect with time element of Stokes equation.
   ８種類の微分幾何では、それぞれの時間の固有値が違う。
   閉3次元多様体上で微分幾何の切除、分解によって時間の性質が決まる。
   This manifold gut theory from described with zeta function to
   catastrophe for non tree of routs result on sergery of space system.
   閉３次元多様体に統合されると、ゼータ関数になり、各幾何に分解された場合
   に、非分岐から、この上、sergeryの結果が決まる。

   各微分構造においての時間発展での熱エネルギーの変化は$E=mc^2, E=m^2 c^2 $
   のthree manifoldが微分幾何構造の変化にそれぞれ対応する。
   Each geometry point interacte with exchange of three manifold in
   Seifert structure from special relativity of equation for
   heat energy fluentations on time developyed from this extermaination.

   $$ \nabla ({i \hbar^{\nabla}})^{\oplus L}, \bigoplus ({i \hbar^{\nabla}})^{
	   \oplus L}, \square ({i \hbar^{\nabla}})^{\oplus L} $$
   These operator of equation on summatative manifolds from emerged with
   element of particle conclution, Euler product is resulted from this
   operator expalanation of locality insectations.
   これらの作用素が加減乗除の生成式の元、Euler Productの結論による
   作用素生成の論理素子でもある。
   Moreover, this eight geometry of differential operator are
   constructed with four pattern of Jones manifold from
   summatative formula and this system extate with special relativity
   references. This decieved of elemet on summatative equation routed
   to internext in real and imaginary pole on complex dimension,
   and this dimension explation with Stokes theorem defined too.
   And this defined circutation of Yacobi matrix is represected with
   Knot theory with anstate with Jones manifold.
   その上に、8種類の微分幾何は、Jones多項式の4パターンでの構成される
   差分と加群方程式から、特殊相対性理論をも思わせる、この差分方程式
   $$ {d \over d \gamma}\Gamma = e^{f} + e^{-f} \ge e^{f} - e^{f} $$
   から、8種類の代数幾何が、ストークスの定理と同じく、このJones多項式
   の固有周期が、すべての時間の流れの基軸をも内包してもいることが、
   わかる。このJones多項式は、結び目の理論をも、この周期が言い表してもいる。
   $$ {d \over d \gamma}\Gamma + \int C dx_m = 2(\cos(i x \log x) - i \sin
   (i x \log x)) $$
    $$ e^{- \theta} = \cos(i x \log x) - i \sin(i x \log x)$$
    $$ \int \Gamma(\gamma)^{'}dx_m = 2(\cos (i x \log x) 
    - i \sin(i x \log x)) $$

 カルーツァ・クライン空間の方程式は、
    $$ ds^2 = g_{\mu\nu}(x)dx^{\mu}dx^{\nu} + \psi^2(x)(dx^2 + \kappa^2
    A_{\mu}(x)dx^{\nu})^2 $$
    この式は、
    $$ ds^2 = -N(r)^2dt^2 + \psi^2(x)(dr^2 + r^2d\theta^2) $$
    $$ ds^2 = -dt^2 + r^{-8\pi Gm}(dr^2 + r^2d\theta^2) $$
    とまとまり、
    $$ dx^2 = g_{\mu\nu}(x)(g_{\mu\nu}(x)dx^2 - dxg_{\mu\nu}(x)) $$
    $$ dx = (g_{\mu\nu}(x)^2dx^2 - g_{\mu\nu}(x)dx g_{\mu\nu}(x))^{1 \over 2} $$
    と反重力と正規部分群の経路和となり、　
    $$ \pi(\chi, x) = i \pi (\chi, x)f(x) - f(x)  \pi (\chi, x) $$
    と基本群にまとまる。
    シュワルツシルト半径は、
    $$ ds^2 = -(1 - {r_s \over r})c^2 dt^2 + {1 \over {{1 - {r_s \over r}}}} 
    dr^2 + r^2 d \theta^2 + r^2 \sin^2 \theta d \psi^2 $$
    これは、カルーツァ・クライン空間と同型となる。
    一般相対性理論は、
    $$ R^{\mu\nu} + {1 \over 2}g_{ij}\Lambda = \kappa T^{\mu\nu} $$
    この多様体積分は、
    $$ \int \kappa T^{\mu\nu}\mathrm{dvol} = \int  \left(R^{\mu\nu} + 
    {1 \over 2}g_{ij}\Lambda\right)\mathrm{dvol}  $$
    ガンマ関数のおける大域的微分多様体は、これと同型により、
    $$ \int \Gamma \cdot {d \over {d \gamma}}{\Gamma dx_m} \le \int \Gamma
    dx_m + {d \over d \gamma}\Gamma $$
    $$ = \int \Gamma(\gamma)^{'}dx_m $$
    オイラーの定数の多様体積分は、
    $$ \int \left({{\int {1 \over x^s}}dx - \log x} \right)\mathrm{dvol} $$
    この解は、シュワルツシルト半径と同型より、
    $$ = e^f - e^{-f} \le {e^f + e^{-f}} $$
    次元の単位は多様体より、大域的微分多様体とオイラーの定数
    の多様体積分の加群分解は、オイラーの公式の三角関数の虚数の度解より、
    $$ {d \over df}F + \int C dx_m = 2 (\cos (i x \log x) 
    - i \sin (i x \log x)) $$
    すべては、大域的微分多様体の重力と反重力方程式に行き着き、
    $$ {d \over df}F \ge {d \over df}\int \int{1 \over (x \log x)^2}dx_m
    + {d \over df}\int \int{1 \over (y \log y)^{1 \over 2}}dy_m $$
    と求まる。


  ８種類の微分幾何での、それぞれの時間の固有値の動静は、
  ストークスの流体力学での、熱エネルギーの各流れの仕方で表されている。
  種数１の2種が種数０の球体に両辺で交わっていると、
一般相対性理論における重力理論は大域的微分多様体と積分多様体についての 単体量を求めるための空間の対数関数における不変性を記述するプランクスケール と異次元の宇宙におけるスケール、ウィークスケールと言われているanti-D-brane として、この２種類の計量をガンマ関数とベータ関数としてオイラーの定数 とそれによる連分数が微分幾何の量子化と数式の値を求めると同じという 予知と推測値から、広中平祐定理の４重帰納法のオイラーの公式による 多様体積分と同類値として、ペレルマン多様体がサーストン空間に成り立つ と同じく、広中平祐定理がヒルベルト空間で成り立つとしての、 これら２種の多様体の場の理論が、ゼータ関数がAdS5多様体で成り立つことと、 不変量として、ゼータ関数がこの3種の多様体の場の理論のバランスをとる
  理論として言えることである。
$$ ||ds^2|| = \lim_{x \to \infty} [\delta(x) \int \int \int \pi \left( \sum_{k=0}^{\infty}
		{{}^n\sqrt{p},x \over n} \right)
^{1 \over 2}d \tau]^{\mu\nu} $$
  This norm is component of fields in Hilbert manifold of space theorem rebuilt with
Mebius space in Gamma Function on Beta function fill into power of rout fundamental
group.

  　 Gravity of general relativity theory describe with Cutting of space
in being discatastrophed from Global differential and integral manifold
of scaled levetivity in plank scal and the other vector of universe
scale, these two scale inspectivity of Gamma and Beta function escourted
with these manifold experted in result of Differentail geometry of
quantum level manifold equal with Euler product and continue parameter,
moreover this product is relativity of Hironaka theorem in Four assembled
integral of Euler equation.

   \newcommand{\variint}{%
	\begin{picture}(15,30)
		\put(0,0){
			$ \iint $}
		\put(0,5){\line(15,0){15}}
	\end{picture} %
}

\newcommand{\variiint}{%
	\begin{picture}(15,15)
		\put(0,0){
			$ \iiint $}
		\put(0,5){\line(15,0){20}}
	\end{picture} %
}



 \newcommand{\alearletter}{%
	 \begin{picture}(10,20)
		 \put(0,0){
			 $ \square $}
		 \put(0,0){\line(1,1){10}}
	 \end{picture} %
	 }

 重力場理論の式は、
  Gravity equation is
  $$ \square = - {16 \pi G \over c^4}T^{\mu\nu} $$
  This equation quated with being logment of formula, and
  this formula divided with universe of number in prime zone,
  therefore, this dicided with varint equation is monotonicity of
  being composited with Weil's theorem united for Gamma function.
  この式は、対数を宇宙における数により求める素数分布論として、この
  大域的積分分断多様体がガンマ関数をヴェイユ予想を根幹とする単体量
  として決まることに起因する商代数として導かれる。
  $$ \scalebox{1}[2]{\alearletter} = 
  {{{{8 \pi G} \over {c^4}}T^{\mu\nu}}/ \log x} $$
  $$ {}^{t}\scalebox{1}[2]{\variiint}\mathrm{cohom D_{\chi}[I_m]}$$
  $$ = \oint {(px^n + qx + r)^{\nabla l}} $$
  $$ {d \over d l}L(x,y) = 2 \int ||\sin 2x||^{2}d \tau $$
  $$ {d \over d \gamma}\Gamma $$
  この関数は大域的微分多様体としてのアカシックレコードの合流地点として、
  タプルスペースを形成している。
  This function esterminate with acasic record of global differential
  manifolds.
  $$ = [i \pi (\chi, x), f(x)] $$
  それにより、この多様体は基本群をアカシックレコードの相対性としての
  存在論の実存主義から統合される多様体自身としてのタプルスペースの
  池になっている。
  And this manifold from fundemental group esterminate with
  also this manifold estimate relativity of acasic record.
  $$ {d \over d \gamma}\Gamma = \int C dx_m = \int(\int {1 \over x^s}dx
  - \log x)\mathrm{dvol} $$
  また、このアカシックレコードはオイラーの定数のラムダドライバー
  にもなっている。
  More also this record tupled with lake of Euler product.
  $$ \bigoplus({i \hbar^{\nabla}})^{\oplus L} $$
  $$ x^{{1 \over 2} + iy} = e^{x \log x} $$
  それゆえに、この定数はゼータ関数と微分幾何の量子化を因数にもつ
  素因子分解の式にもなっている。
  Therefore, this product also constructed with differential geometry
  of quantum level equation and zeta function.
  $$ = C $$ 
  そして、この関数はオイラーの定数から広中平祐定理による4重帰納法の
  オイラーの公式からの多様体積分へとのサーストン空間のスペクトラム関数
  ともなっている。
  And this function of Euler product respectrum of focus with
  Heisuke Hironaka manifold in four assembled of integral Euler equation.. 

  $$ C =  b_0 + \cfrac{c_1}{b_1 +
	  \cfrac{c_2}{b_2 +
	  \cfrac{c_3}{b_3 +
	  \cfrac{c_4}{b_4 + \cdots}}}} $$
  この方程式は指数による連分数としての役割も担っている。  
  This equation demanded with continued number of step function.
  $$ x^{{1 \over 2} + iy} = e^{x \log x} $$
  $$ (2.71828)^{2.828} = a^{a + b^{b + c^{c + d^{d \cdots }}}} $$
  $$ = \int e^f \cdot x^{1 - t}dx $$
  This represent is Gamma function in Euler product.
  Therefore this product is zeta function of global differential equation.

  $$ {d \over d f}F(v_{ij},h) = \int {e^{-f}\mathrm{dV}}
  [- \Delta v + \nabla_{i}\nabla_{j}
  v_{ij} - R_{ij}v_{ij} - v_{ij}\nabla_{i}\nabla_{j} + 2 <\nabla f,\nabla h>
  + (R + \nabla f^2)({v \over 2} - h)]$$
  $$ = {d \over d f}F = {2 {\int (R + \nabla_{i}\nabla_{j} f)^2} \over 
  { - (R + \Delta f)}}\mathrm{dm} $$
  これらの方程式は8種類の微分幾何の次元多様体として、そして、これらの
  多様体は曲平面による双対性をも生成している。 
  そして、このガウスの曲平面は、大域的微分多様体と微分幾何の量子化から
  素因数を形成してもいる。
  These equation are eight differential geometry of dimension calyement.
  And these calyement equation excluded into pair of dimension surface.
  This surface of Gauss function are global differential manifold,
  and differential geometry of quantum level.
  $$ F \ge {d \over d f}\int \int{1 \over (x \log x)^2}dx_m 
  + {d \over d f}\int \int{1 \over (y \log y)^{1 \over 2}}dy_m $$

  微分幾何の量子化はオイラーの定数とガンマ関数が指数による連分数
  としての不変性として素因数を形成していて、このガウスの曲平面
  による量子力学における重力場理論は、ダランベルシアンの切断多様体
  がこの大域的切断多様体を付加してもいる。
  Differential geometry of quantum level constructed with Euler product
  and Gamma function being discatastrophed from continued fraction style.
  Gravity equation lend with varint of monotonicity of level expresented
  from gravity of letter varient formula.
  これらの方程式は基本群と大域的微分多様体をエスコートしていもいて、
  ヴェイユ予想がこのダランベルシアンの切断方程式たちから
  輸送のポートにもなっている。ベータ関数とガンマ関数がこれらのフォームラ
  の方程式を放出してもいて、結果、これらの方程式は広中平祐定理の複素
  多様体とグリーシャ教授によるペレルマン多様体からサポートされてもいる。
  この2名の教授は、一つは抽象理論をもう一方は具象理論を説明としている。
  These equation escourted into Global differential manifold and
  fundemental equation.
  Weil's theorem is imported from this equation in gravity of letters.
  Beta function and Gamma function are excluded with these formulas.
  These equation comontius from Heisuke Hironaga of complex manifold
  and Gresha professor of Perelman manifold.
  These two professos are one of abstract theorem and the other of
  visual manifold theorem.

\section{Hilbert manifold in Mebius space \\
        this element of Zeta function on integrate of fields}

 Hilbert manifold equals with Von Neumann manifold, and this fields is concluded with Glassman manifold,
the manifold is dualty of twister into created of surface built.
These fields is used for relativity theorem and quantum physics.
$$ ||ds^2|| = \lim_{x \to \infty} [\delta(x) \int \int \int \pi \left( \sum_{k=0}^{\infty}
		{{}^n\sqrt{p},x \over n} \right)
^{1 \over 2}d \tau]^{\mu\nu} $$
  This norm is component of fields in Hilbert manifold of space theorem rebuilt with
Mebius space in Gamma Function on Beta function fill into power of rout fundamental
group.
And this result with $AdS_5$ space time in Quantum caos of Minkowsky of manifold in
abel manifold.
Gravity of metric in non-commutative equation is Global differential equation
conbuilt with Kaluza-Klein space. Therefore this mechanism is $T^{\mu\nu}$ tensor
is equal with $R^{\mu\nu}$ tensor. And This moreover inspect with laplace operator
in stimulate with sign of differential operator.
Minus of zone in Add position of manifold is Volume of laplace equation rebuilt with
Gamma function equal with summative of manifold in Global differential equation,
this result with setminus of zone of add summative of manifold, and construct with
locality theorem straight with fundamental group in world line of surface, this
power is boson and fermison of cone in hyper function.

$$ V(\tau) = [f(x),g(x)] \times [f^{-1}(x), h(x)] $$
$$ \Gamma(p,q) = \int e^{-x} x^{1 - t}dx $$
$$ = \beta(p,q) $$
$$ = \pi (f(\chi,x),x) $$
$$ ||ds^2|| = \mathcal{O}(x)[(f(x) \circ g(x))^{\mu\nu}]dx^{\mu}dx^{\nu} $$
$$ = \lim_{x \to \infty} \sum_{k=0}^{\infty}a_k f^k $$
$$ G^{\mu\nu} = {\partial \over \partial f}\int [f(x)^{\mu\nu} \circ G(x)^{\mu\nu}dx^{\mu}dx^{\nu}]^{\mu\nu}dm $$
$$ = g_{\mu\nu}(x)dx^{\mu}dx^{\nu} - f(x)^{\mu\nu}dx^{\mu}dx^{\nu} $$
$$ [i \pi(\chi,x),f(x)] = i \pi f(x) - f(x) \pi(\chi,x) $$
$$ T^{\mu\nu} = (\lim_{x \to \infty} \sum_{k=0}^{\infty}\int \int [V(\tau) \circ S^{\mu\nu}(\chi,x)]dm)
^{\mu\nu}dx^{\mu}dx^{\nu} $$
$$ G^{\mu\nu} = R^{\mu\nu}T^{\mu\nu} $$
$$ \begin{vmatrix} D^m & dx \\
dx & \sigma^m \end{vmatrix} \begin{vmatrix} \cos \theta & - \sin \theta \\
\sin \theta & \cos \theta \end{vmatrix} \begin{vmatrix} x \\
y \end{vmatrix} = \begin{pmatrix} 1 & 0 \\ 
0 & - 1 \end{pmatrix}^{1 \over 2} $$
$$ \sigma^m \begin{bmatrix} \delta(x) & - 1 \\
1 & \epsilon(x) \end{bmatrix}^{1 \over 2} = \begin{pmatrix} i & 0 \\
0 & - i \end{pmatrix} $$
$$ V(M) = {\partial \over \partial f} ({}^N \int [f \smallsetminus M]^{\oplus N})^{\mu\nu} dx^{\mu}dx^{\nu} $$
$$ V(M) = \pi(2 \int \sin^2 dx)\oplus {d \over df}F^M dx_m $$
$$ \lim_{x \to \infty}\sum_{k=0}^{\infty} a_k f^k = \int (F(V)dx_m)^{\mu\nu}dx^{\mu}dx^{\nu} $$
$$ \bigoplus_{k=0}^{\infty}[ f \smallsetminus g] = \vee(M \wedge N) $$
$$ \pi_1 (M) = e^{-f 2 \int \sin^2 x dm} + O(N^{-1}) $$
$$ = [i \pi(\chi,x),f(x)] $$
$$ M \circ f(x) = e^{-f \int \sin x \cos x dx_m} + \log (O(N^{-1})) $$

 Non-Symmetry space time.
 $$ {d \over dL}V(\tau) = {d \over df}\int \int_M {1 \over (x \log x)^2}dx_m + {d \over df}\int \int_M
 {1 \over (y \log y)^{1 \over 2}}dy_m $$
 $$ \epsilon S(\nu) = \square_v \cdot {\partial \over \partial \chi}({}^5 \sqrt{\wedge g^2})d\chi $$
 Differential Volume in $AdS_5$ graviton of fundamental rout of group.
 $$ \wedge (F^m_t)^{''} = {1 \over 12}g_{ij}^2$$
 Quarks of other dimension.
 $$ \pi(V_{\tau}) = {e^{-(\sqrt{\pi \over 16}\log x)}}^{\delta} \times {1 \over (x \log x)} $$
 Universe of rout, Volume in expanding space time.
 $$ {d \over dt}(g_{ij})^2 = {1 \over 24}(F^m_t)^2 $$
 $$ m^2 = 2 \pi T \left({26 - D_n \over 24}\right) $$
 This quarks of mass in relativity theorem, and fourth of universe in three manifold of one dimension surface,
 and also this integrate of dimension in conbuilt of quarks.
 $$ g_{ij} \wedge \pi(\nu_{\tau}) = e^{-2 \pi T |\psi|}[\eta_{\mu\nu} + \bar{h}_{\mu\nu}(x)]dx^{\mu}dx^{\nu} 
 + T^2 d \psi^2 $$
 Out of rout in $AdS_5$ space time.
 $$ ||ds^2|| = g_{ij}\wedge \pi(\nu_v) $$
 $AdS_5$ norm is fourth of universe of power in three manifold out of rout.


  Sphere orbital cube is on the right of hartshorne conjecture by the right equation,
this integrate with fourth of piece on universe series, and this estimate with one of three manifold.
Also this is blackhole on category of symmetry of particles.
These equation conbuilt with planck volume and out of universe is phologram field, this field construct with
electric field and magnitic field stand with weak electric theorem.
This theorem is equals with abel manifold and $AdS_5$ space time.
Moreover this field is anti-brane and brane emerge with gravity and antigravity equation restructed with.
Zeta function is existed with Re pole of ${1 \over 2}$ constance reluctances.
This pole of constance is also existed with singularity of complex fields.
Von Neumann manifold of field is also this means.
$$ ||ds^2|| = \lim_{x \to \infty} [\delta(x)\int \int \int \pi \left( \sum_{k=0}^{\infty}{{}^n\sqrt{p},x \over n}\right)
^{1 \over 2}d \tau]^{\mu\nu} $$
$$ e^{i \theta} = \cos \theta + i \sin \theta $$
$$ e^{x \log x} = x^{{1 \over 2} + iy}, x \log x = \log (\cos \theta + i \sin \theta) $$
$$ = \log \cos \theta + i \log \sin \theta $$
This equation is developed with frobenius theorem activate with logment equation.
This theorem used to
$$ \log(\sin \theta + i \cos \theta)= \log(\sin \theta - i \cos \theta) $$
$$ \log \left({\sin \theta \over i \cos \theta}\right) = -2 R_{ij}, {d \over dt}g_{ij}(t) =  - 2 R_{ij} $$
$$ \mathcal{O}(x) = {\zeta(s) \over \sum\limits_{k=0}^{\infty}a_k f^k} $$
$$ \mathrm{Im} f = \mathrm{ker} f, \chi(x) = {\mathrm{ker}f \over \mathrm{Im}f} $$ $$ H(3) = 2, \nabla H(x) = 2, 
 \pi(x) = 0 $$
 This equation is world line, and this equation is non-integrate with relativity theorem of rout constance.

$$ [f(x)] = \infty, ||ds^2|| = \mathcal{O}(x)[\eta_{\mu\nu} + \bar{h}_{\mu\nu}(x)]dx^{\mu}dx^{\nu} + T^2 d^2 \psi $$
$$ T^2 d^2 \psi = [f(x)] , T^2 d^2 \psi = \lim_{x \to 1}\sum_{k=0}^{\infty}a_k f^k $$

 These equation is equals with M theorem. For this formula with zeta function into one of universe in four dimensions.
Sphere orbital cube is on the surface into $AdS_5$ space time, and this space time is Re pole and Im pole of 
constance reluctances.


  Gauss function is equals with Abel manifold and Seifert manifold. Moreover this function is infinite time,
and this energy is cover with finite of Abel manifold and Other dimension of seifert manifold on the surface.


  Dilaton and fifth dimension of seifert manifold emerge with quarks and Maxwell equation, this power
is from dimension to flow into energy of boson. 
This boson equals with Abel manifold and $AdS_5$ space time, and this space is created with Gauss function.


  Relativity theorem is composited with infinite on D-brane and finite on anti-D-brane, This space time
restructed with element of zeta function.
Between finite and infinite of dimension belong to space time system, estrald of space element have with 
infinite oneselves. Anti-D-brane have infinite themselves, and this comontend with fifth dimension of AdS5 
have for seifert manifold, this asperal manifold is non-move element of anti-D-brane and move element 
of D-brane satisfid with fifth dimension of seifert algebra liner.
Gauss function is own of this element in infinite element.
And also this function is abel manifold.
Infinite is coverd with finite dimension in fifth dimension of AdS5 space time.
Relativity theorem is this system of circustance nature equation.
AdS5 space time is out of time system and this system is belong to infinite mercy.
This hiercyent is endrol of memolite with geniue of element.
Genie have live of telomea endore in gravity accesorlity result.
AdS5 space time out of over this element begin with infinite assentance.
Every element acknowlege is imaginary equation before mass and spiritual envy.

  $$ T^2 d^2 \psi = \lim_{x \to 1}\sum_{k=0}^{\infty}a_k f^k $$
  $$ \lim_{x \to 1}\sum_{k=0}^{\infty}a_k f^k = [T^2 d^2 \psi] $$

  $$ {d \over dL}V(\tau)= {d \over df}\int \int_M({}^5\sqrt{x^2})d\Lambda 
  + {d \over df}\int \int_M {}^N({}^3\sqrt{x})^{\oplus N}d\Lambda $$
  $$ {}^M(\vee(\wedge f \circ g)^N)^{1 \over 2} = {d \over df}\int \int_M{1 \over (x \log x)^2}dx_m
  + {d \over df}\int \int_M{1 \over (y \log y)^{1 \over 2}}dy_m $$
  $$ ||ds^2|| = \mathcal{O}(x)[\eta_{\mu\nu} + \bar{h}_{\mu\nu}(x)]dx^{\mu}dx^{\nu}
  + T^2 d^2 \psi $$
  $$ \mathcal{O}(x) = e^{-2\pi T|\psi|} $$
  $$ G^{\mu\nu} = R_{\mu\nu} T^{\mu\nu} $$
  $$ = - {1 \over 2}\Lambda g_{ij}(x) + T^{\mu\nu} $$

  Fifth dimension of eight differential structure is integrate with one geometry element,
Four of universe also integrate to one universe, and this universe represented with symmetry formula.
Fifth universe also is represented with seifert manifold peices.
All after universe is constructed with six element of circuatance.
This aquire maniculate with quarks of being esperaled belong to.

\section{Imaginary equation in AdS5 space time create with \\
       dimension of symmetry}

  D-brane and anti-D-brane is composited with all of series universe emerged for one geometry of 
dimension, this gravity of power from D-brane and anti-D-brane emelite with ancestor.
Seifert manifold is on the ground of blackhole in whitehole of power pond of senseivility.
Six of element of quarks and universe of pieces is supersymmetry of mechanism resolved with

hyper symmetry of quarks constructed to emerge with darkmatter.
This darkmatter emerged with big-ban of heircyent in circumstance of phenomena.

  D-brane and anti-D-brane equations is
  $$ {d \over dL}V(\tau) = {d \over df}\int \int_M {1 \over (x \log x)^2}dx_m
  + {d \over df}\int \int_M {1 \over (y \log y)^{1 \over 2}}dy_m $$
  $$ C = \int \int {1 \over (x \log x)^2}dx_m $$

  Euler constance is quantum group theorem rebuild with projective space involved with.


   虚数方程式は、反重力に起因するフーリエ級数の励起を生成する。それは、人工知能を生み出す、
   ５次元時空にも、この虚数方程式は使われる。AdS5の次元空間は、反ド・シッター時空の
   D-braneとanti-D-braneのcomformal場を生み出す。ホログラフィー時空は、この量子起因によるものである。
   2次元曲面によるブラックホールは、ガンマ線バーストによる５次元時空の構造から観測される。
   空間の最小単位によるプランク定数は、宇宙の大域的微分多様体から導き出される、AdS5の次元空間の
   準同型写像を形成している。これは、最小単位から宇宙の大きさを導いている。最大最小の方程式は、
   相加相乗平均を形成している。時間と空間は、宇宙が生成したときから、宇宙の始まりと終わりを既に生み出している。
   宇宙と異次元から、ブラックホールとホワイトホールの力がわかり、反重力を見つけられる。オイラーの定数は、
   この量子定数からわかる。虚数の仕組みはこの量子スピンの産物である。
   オイラーの定数は、 この虚時空の斥力の現存である。それは、非対称性理論から導かれる。
   不確定性原理は、AdS5のブラックホールとホワイトホールを閉３次元多様体に統合する５次元時空から求められる。
   位置と運動エネルギーが、空間の最小単位であるプランク定数を宇宙全体にする微細構造定数からわかり、
   面積確定から、アーベル多様体を母関数に極限値として、ゼータ関数をこの母関数に不変式として、
   ２種類ずつにまとめる４種類の宇宙を形成する８種類のサーストンの幾何化予想から導き出される。
   この閉３次元多様体は、ミラー対称性を軸として、６種類の次元空間を一種類の異次元宇宙と同質ともしている。
   複素多様体による特異点解消理論は、この原理から求められる。この特異点解消理論は、２次元曲面を３次元多様体
   に展開していく、時空から生成される重力の密度を反重力と等しくしていく時間空間の４次元多様体と虚時空
   から求められる。
   ヒルベルト空間は、フォン・ノイマン多様体とグラスマン多様体をこのサーストンの幾何化予想を場の理論
   既定値として形成される。
   この空間は、ミンコフスキー時空とアーベル多様体全体を表している。そして、この空間は、球対称性を
   複素多様体を起点として、大域的トポロジーから、偏微分を作用素微分として時空間をカオスから
   ずらすと５次元多様体として成立している。
   これらより、３次元多様体に２次元射影空間が異次元空間として、AdS5空間を形成される。
   偏微分、全微分、線形微分、常微分、多重微分、部分積分、置換微分、大域的微分、単体分割、
   双対分割、同調、ホモロジー単体、コホモロジー単体、群論、基本群、複体、マイヤー・ビーートリス完全系列、
   ファン・カンペンの定理、層の理論、コホモルティズム単体、CW複体、ハウスドロフ空間、線形空間、位相空間、
   微分幾何構造、モースの定理、カタストロフの空間、ゼータ関数系列、球対称性理論、スピン幾何、
   ツイスター理論、双対被覆、多重連結空間、プランク定数、フォン・ノイマン多様体、グラスマン多様体、
   ヒルベルト空間、一般相対性理論、反ド・シッター時空、ラムダ項、D-brane,anti-D-brane,コンフォーマル場、
   ホログラム空間、ストリング理論、収率による商代数、ニュートンポテンシャルエネルギー、剛体力学、
   統計力学、熱力学、量子スピン、半導体、超伝導、ホイーストン・ブリッチ回路、非可換確率論、Connes理論、
   これら、演算子代数を形成している、微分・積分作用素が、ヒルベルト空間に存在している多様体の特質を
   全面に押し出して、いろいろな多様体と関数そして、群論を形作っている。


$$ \begin{vmatrix} D^m & dx \\
dx & \sigma^m \end{vmatrix} \begin{vmatrix} \cos \theta & - \sin \theta \\
\sin \theta & \cos \theta \end{vmatrix} \begin{vmatrix} x \\
y \end{vmatrix} = \begin{pmatrix} 1 & 0 \\ 
0 & - 1 \end{pmatrix}^{1 \over 2} $$
$$ \sigma^m \begin{bmatrix} \delta(x) & - 1 \\
1 & \epsilon(x) \end{bmatrix}^{1 \over 2} = \begin{pmatrix} i & 0 \\
0 & - i \end{pmatrix} $$
$$ (D^m,dx)\cdot (\cos \theta, \sin \theta) \times (dx, \partial^m)\cdot (\cos \theta, \sin \theta) $$

$$ \begin{pmatrix} 1 && 0 \\
   0 && i \end{pmatrix} \beta(x,\theta) \begin{pmatrix} x \\
   y \end{pmatrix} $$
   $$ = \begin{pmatrix} i && 0 \\
   0 && - 1 \end{pmatrix} $$
   $$ l = {\sqrt{\hbar G \over c^3}}, \sigma^m \cdot \begin{pmatrix} \delta(x) && - 1 \\
   1 && \epsilon (x) \end{pmatrix}^{1 \over 2} = \begin{pmatrix} i && 0 \\
   0 && - i \end{pmatrix} $$
   String theorem also imaginary equation of pole with the fields of manifold.
   $$ \left({\partial \over \partial \tau}f(x,y,z)\right)^{3'} = A^{\mu\nu} $$
   $$ {d \over dt}g_{ij}(t) = - 2 R_{ij} $$
   Rich equation is all of top formula in integrate theorem.
      $$ = (D^m,dx) \cdot (\cos \theta, \sin \theta) \times (dx, \partial^m)\cdot (\cos \theta, \sin \theta) $$
      Dalanverle equation equals with abel manifold.






\title{慣性の法則と、回転体による反重力 \\
タイミングとともに、別ベクトルとして重力を加える法則}
\author{Masaaki Yamaguchi}


   慣性の法則の慣性力が働くときと同時に、
   前もって、タイミングを計り、慣性力とともに、
   逆向きのベクトルとして、慣性力の力と同じエネルギーの
   重力を生成して、慣性力の影響を無くす。この力を反重力という。
   慣性力と逆向きのベクトルとして、反重力として、
   慣性力に加えて、別のベクトルとする。

   反重力の本質は、重力のベクトルの違いであり、
   核エネルギーが特殊相対性理論と一般相対性理論を使って、
   放射性物質吸収体に、この原理を使い、核エネルギーをリサイクルしている。
   これを宇宙人は知っていて、UFOにこの原理を使っている。

   Higgs場は、重力であり、周りを反重力が覆っている。
   一般相対性理論のベクトルの向きを変えた力の装置が、
   ラムダ・ドライバであり、反重力である。

   重力は熱エネルギーである。ベクトルの向きが違う。
   反重力は吸収体である。
   慣性の法則が働いているときに、これが手に入ったが、離れたになり、
   重力をベクトルの向きを変えて、同じエネルギーを加えた場合、
   量子コンピュータが、そのときの差分を使って、反重力が慣性の法則
   を変えて、宙に浮くUFOにしている。 

   まとめると、慣性の法則のときに、重力のベクトルを違う向きとして、
   加えると、手に入ったが、離れたになり、これが反重力である。 

   UFOが現れる前に、ヘリコプターが現れて、そのあとに、UFOが実際に現れる。
   慣性の法則のしばらくして、反重力として、UFOとして、
   宇宙人が人に教えている。

   反重力の生成は、回転体の電磁気力からくるちからであり、ローレンツ力と
   同じでもあり、機体の内部で生成される力であり、そのために、放射性
   物質吸収体を結晶石として、このエネルギーからくる放射力を吸収する
   ために、結晶石に反重力の放射線を吸収する。その上に、外部の
   慣性力を内部の反重力で、手に入ったら、離れたをしている。


\title{反重力発生装置の各構成物質}
\author{Masaaki Yamaguchi}


  原子振動子　セシウム Cs


  形状記憶合金 (Fe・Co60・Pt・Al)$H_2 SO_4, Al_2 O_3$


  反重力発生器 Pr(パラジウム合金)電磁場生成He,$H_2$Mg,Al(摩擦熱)


  結晶石 Co60,Pt,Prの反陽子


  緩衝剤 慣性力感知器 有機化合物とS化合物(シクロアルカン$C_n H_{2n}$
  方位と位置探知機にも使う


  量子コンピュータの量子素子 Pt,Ag,Au（電子と電流のエネルギー経路)半導体に
  使う
  回路の合成演算子 


  $$ x^y = {1 \over {y^x}} \to $$

  $$ \pi(\chi,x) = [i\pi(\chi,x),f(x)] $$
  $$ {d \over df}F = F^{{f}^{'}} $$


  ディスプレイの電磁迷彩 (Al,Mg) $\to$ S合成子 $Si O_2 $
  プリズム C

\end{document}
\documentclass{jsarticle}
\usepackage[dvipdfmx]{graphicx}
\usepackage{chemfig}
\usepackage{amsmath}
\usepackage{amssymb}
\begin{document}
\title{Fiber of layzerium in amalugam and stephany of parm \\
in antigravity from electricity \\
and magnity of non condition's power}
\author{Masaaki Yamaguchi}
\date{}
\maketitle

   大域的多重積分と大域的無限微分多様体を定義すると、
   Global assemble manifold defined with infinity of differential
   fields to determine of definition create from Euler product and
   Gamma function for Beta function being potten from integral manifold
   system. Therefore, this defined from global assemble manifold from
   being Gamma of global equation in topological expalanations.

   $$ \int f(x)dx = \int \Gamma(\gamma)^{'}dx_m $$
   $$ = 2 (\cos (i x \log x) - i \sin(i x \log x)) $$
   $$ \left({{\int f(x) dx} \over \log x}\right) = \lim_{\theta \to \infty}
    \left({{\int f(x) dx} \over \theta}\right) = 0,1 $$
    $$ e^{i \theta} = \cos \theta + i \sin \theta $$
    $$ \left({\int f(x) dx}\right)^{'} = 2 (i \sin(i x \log x) - \cos (i x \log
    x)) $$
    $$ = 2(- \cos(i x \log x) + i \sin(i x \log x)) $$
    $$ (\cos (i x \log x) - i \sin(i x \log x))^{'} $$
    $$ = {d \over {d {e^{i\theta}}}}\left((\cos, - \sin )\cdot
    (\sin, \cos)\right)  $$
    $$ \int \Gamma(\gamma)^{''}dx_m = \left({\int \Gamma(\gamma)^{'}}dx_m
    \right)^{\nabla L} = \left({{\int \Gamma dx_m} \cdot {{d \over d\gamma}
    \Gamma}}\right)^{\nabla L} \le 
    \left({\int \Gamma dx_m + {d \over d\gamma}\Gamma}\right)^{\nabla L} \le 
    \left({e^{f} - e^{-f}} \le {e^{-f} + e^{f}}\right)^{'}  $$
    $$ = 0,1 $$

    Then, the defined of global integral and differential manifold
    from being assembled world lines and partial equation of deprivate
    formula in Homology manifold and gamma function of global topology
    system, this system defined with Shwaltsshild cicle of norm from
    black hole of entropy exsented with result of monotonicity for
    constant of equations.
    以上 より、大域的微分多様体を大域的 2重微分多様体として、
    処理すると、ホモロジー多様体では、種数が１であり、特異点では、種数が０
    と計算されることになる。ガンマ関数の大域的微分多様体では、
    シュバルツシルト半径として計算されうるが、これを大域的2重微分で
    処理すると、ブラックホールの特異点としての解が無になる。

    川島隆太先生が、私みたいな人が、頭の中で偏微分や多重積分、部分積分、
    共変微分や、置換積分や、複素微分や複素積分を考えていると思っていた
    と言っていたが、それ以上の大域的微分や大域的2重微分や、
    大域的多重積分や大域的偏微分や大域的部分積分、微分幾何の量子化、
    単体積分、単体微分などの、大学院生でも考えない、思いもしない、
    実際の式を書くと誰でも、その関数と多様体が存在するなと思うしか
    できない、私の記述式を見ると納得する理論式でもあると、
    私も自慢できる理論であると自負できる。

  $$ \log x|_{g_{ij}}^{\nabla L} = f^{f^{'}}, F^f|_{g_{ij}}^{\nabla L}
  |: x \to y, x^p \to y $$ 
  $$  f(x) = \log x = p \log x, f(y) = p \log x  $$
  大域的偏微分方程式と大域的多重積分は、それぞれ次のように成り立っている。
  It's defined with global partial deprivate formula and assemble manifold. 
  $$ {d^2 \over df^2}F = F^{f^{'}}\cdot f^{f^{''}},  $$
  $$ \int \int F dx_m = F^f \cdot F^{(f)^{'}} $$
  $$ {d \over df dg}(f,g) = (f \cdot g)^{f^{'} + g^{'}} $$
  大域的部分積分も、次のように成り立っている。
  And, global parcial integral equation also established with global
  topology computations. therefore, this circutation exclude for all of
  topology extension ideas.
  $$ (F^f \cdot G^g) = \int (f \cdot g)^{f^{'} + g^{'}} $$
  $$ \int \int F\cdot G dx_m = [F^f \cdot G^g] - 
  \int (f \cdot g)^{f^{'} + g^{'}} $$
  大域的部分積分の計算は、
  $$ \int {d \over {df dg}}FG = \int F^{f^{'}}G + \int F G^{g^{'}} $$
  $$ \int F^{f^{'}}Gdx_m = [F^f G^g] - \int F G^{g^{'}}dx_m$$
  大域的商代数の計算は、
  Global quato algebra computations also,
  $$ \left({F(x) \over G(x)}\right)^{(fg)^{'}} = {F^{f^{'}}G - FG^{g^{'}}
  \over G^{g}} $$
  大域的偏微分方程式は、縮約記号を使うと、

  Global partial manifold explain with reduction of operator for being
  represented from partial equation of references.
  $$ {d \over df dg}FG = {d \over df_m}FG $$
  $$ {\partial \over \partial f_m}FG = F^{f^{\mu\nu}}\cdot G
  + F \cdot G^{g^{\mu\nu}}, \int F dx_m = F^f $$
  多様体による大域的微分と大域的積分が、エントロピー式で統一的に表せられる。

  Entorpy computation from integrate for being manifold of global deprivate
  and integrate formulas.
   ベータ関数をガンマ関数へと渡すと、
   Gamma function sented for Beta function of inspiration with monitonicity
   of component with differential and integral being definitions.
   $$ \beta(p,q) = {{\Gamma(p)\Gamma(q)} \over {\Gamma(p + q)}} $$
   ここで、単体的微分と単体的積分を定義すると
   This area defined with different and integrate of component with
   global topology.
   $$ t = \Gamma(x) $$
   $$ T = \int \Gamma(x) dx $$
   これを使うのに、ベータ関数を単体的微分へと渡すために、
   for this system used for being is conbiniate 
   with beta function for componenet of deprivate eqaution.
   $$ \beta(p,q) = - \int {1 \over {t^2}}dt $$
   大域的微分変数へとするが、
   This point be for global differential variable exchanged,
   $$ T_m = \int \Gamma(x) dx_m $$
   この単体的積分を常微分へとやり直すと、
   This also point be for being retried from ordinary differential computation
   for component of integral manifold.
   $$ T^{'} = {t^{'} \over {{\log t}}}dt + C(\text{Cは積分定数}) $$
   これが、大域的微分多様体の微分変数と証明するためには、
   This exceed of proof being for being defined with deprivate variable of
   global differential formula, this successed for true is,
   $$ dx_m = {1 \over {\log x}}dx, dx_m = {(\log x^{-1})}^{'}  $$
   これより、単体的微分が大域的部分積分へと置換できて、
   Therefore, this exchanged from mononotocity deprivation from global
   parital integral manifold is able to, 
   $$ T^{'} = \int \Gamma(\gamma)^{'} dx_m $$

   微分幾何へと大域的積分と大域的微分多様体が、加群分解の共変微分へと
   書き換えられる。
   After all, these exchanged of monotonicity deprivation successed from
   being catastrophe of summativate of partial and assemble of deprivations
   for differential geometory of quantum level to global integral and
   differential manifold.
   $$ \bigoplus T^{\nabla}  = \int T dx_m, \delta(t) = t dx_m  $$


   this exluded of being conclution are which beta function
   evaluate with mononicity from ordinary differential equation
   be resulted from component of deprivation and integral expalanations.
   This cirtutation be resenbled to define with global topoloty of
   extention extern of deprivate and integral of manifolds
   estourced with quantum level of differential geometry be proofed with
   all of equation anbrabed from Euler product.
   これらをまとめると、ベータ関数を単体的関数として、
   これを常微分へと解を導くと、単体的微分と単体的積分
   へと定義できて、これより、大域的微分多様体と大域的積分多様体
   へとこの関数が存在していることを証明できるのを上の式たちから
   わかる。この単体的微分と単体的積分から微分幾何が書けることがわかる。



   虚数の虚数倍した値が超越数のx倍と同じとすると、
   超越数の$2 \pi m$倍が$n=1$でiとなると定義すると、
   次の式たちが導かれる。　
   Imaginary pole circlate with twigled of pole in step function
   from Naipia number of assembled from equalation of defined are
   escourted to be defined with next equations.


   These defined equation are climbate with idea of equation from
   Caltan of imaginary number of circulation.
   カルタンを超えているアイデアと数式たちでもある。

   $$ i^i = (\sqrt{i})^{\sqrt{i}} = e^{x \log x} $$
   $$ e^x = i, e^{2\pi m} = i^n, {d \over d x}e^{2\pi m} = i^n $$
   $$ e^{i \theta} = \cos \theta + i \sin \theta, e^{2 \pi m} = i^n $$
   $$ 2 \pi m = n, e = i \text{($e=i$となる$e^{i \theta} = \cos \theta 
   + i \sin \theta $
   の範囲で$2\pi m$ がある。)}$$
   $$ {d \over d i}i = i^i = e^{x \log x}, m={1 \over 2 \pi},l=2\pi r, 
   \pi = {l \over 2r} $$
   $$ ||ds^2|| = e^{-2 \pi T||\psi||}[\eta_{\mu\nu} + \bar{h}(x)]
   dx^{\mu}dx^{\nu} + T^2 d^2 \psi $$
   $$ e^x = i, e^{2\pi m} = i^n, e = i $$
   $$ {{[\eta_{\mu\nu} + \bar{h}(x)]dx^{\mu}dx^{\nu}}/i} = H \Psi
   = i \hbar \psi, H \Psi = {1 \over i}[H, \Psi] $$
   ハイゼンベルク方程式が$ AdS_5 $多様体の原子レベルの方程式も表せられる。 
   微分幾何の量子化の式は、
   Hisenbelg equation are represeted from AdS5 manifold to particle level
   equalation of quantum levle of differential geometry entranced.
   $$ \bigoplus {({i \hbar} ^{\nabla})}^{\oplus L} = e^{{i x \log x}^{
	   {e^{i x \log x}}^{'}}} $$
   $$ \Gamma = \int e^{-x}x^{1 - t}dx $$
   $$ \gamma = \int e^{-x}x^{1 - t}\log x dx $$
   $$ = \int \Gamma(\gamma)^{'} dx_m$$
   $$ = {d \over d \gamma}\Gamma $$

   Eight differential geometry are each intersect with own level of
   concept from expalanation of Euler product system.
   This component of three manfold sergeried with geometry of destroy and
   desect with time element of Stokes equation.
   ８種類の微分幾何では、それぞれの時間の固有値が違う。
   閉3次元多様体上で微分幾何の切除、分解によって時間の性質が決まる。
   This manifold gut theory from described with zeta function to
   catastrophe for non tree of routs result on sergery of space system.
   閉３次元多様体に統合されると、ゼータ関数になり、各幾何に分解された場合
   に、非分岐から、この上、sergeryの結果が決まる。

   各微分構造においての時間発展での熱エネルギーの変化は$E=mc^2, E=m^2 c^2 $
   のthree manifoldが微分幾何構造の変化にそれぞれ対応する。
   Each geometry point interacte with exchange of three manifold in
   Seifert structure from special relativity of equation for
   heat energy fluentations on time developyed from this extermaination.

   $$ \nabla ({i \hbar^{\nabla}})^{\oplus L}, \bigoplus ({i \hbar^{\nabla}})^{
	   \oplus L}, \square ({i \hbar^{\nabla}})^{\oplus L} $$
   These operator of equation on summatative manifolds from emerged with
   element of particle conclution, Euler product is resulted from this
   operator expalanation of locality insectations.
   これらの作用素が加減乗除の生成式の元、Euler Productの結論による
   作用素生成の論理素子でもある。
   Moreover, this eight geometry of differential operator are
   constructed with four pattern of Jones manifold from
   summatative formula and this system extate with special relativity
   references. This decieved of elemet on summatative equation routed
   to internext in real and imaginary pole on complex dimension,
   and this dimension explation with Stokes theorem defined too.
   And this defined circutation of Yacobi matrix is represected with
   Knot theory with anstate with Jones manifold.
   その上に、8種類の微分幾何は、Jones多項式の4パターンでの構成される
   差分と加群方程式から、特殊相対性理論をも思わせる、この差分方程式
   $$ {d \over d \gamma}\Gamma = e^{f} + e^{-f} \ge e^{f} - e^{f} $$
   から、8種類の代数幾何が、ストークスの定理と同じく、このJones多項式
   の固有周期が、すべての時間の流れの基軸をも内包してもいることが、
   わかる。このJones多項式は、結び目の理論をも、この周期が言い表してもいる。
   $$ {d \over d \gamma}\Gamma + \int C dx_m = 2(\cos(i x \log x) - i \sin
   (i x \log x)) $$
    $$ e^{- \theta} = \cos(i x \log x) - i \sin(i x \log x)$$
    $$ \int \Gamma(\gamma)^{'}dx_m = 2(\cos (i x \log x) 
    - i \sin(i x \log x)) $$

 カルーツァ・クライン空間の方程式は、
    $$ ds^2 = g_{\mu\nu}(x)dx^{\mu}dx^{\nu} + \psi^2(x)(dx^2 + \kappa^2
    A_{\mu}(x)dx^{\nu})^2 $$
    この式は、
    $$ ds^2 = -N(r)^2dt^2 + \psi^2(x)(dr^2 + r^2d\theta^2) $$
    $$ ds^2 = -dt^2 + r^{-8\pi Gm}(dr^2 + r^2d\theta^2) $$
    とまとまり、
    $$ dx^2 = g_{\mu\nu}(x)(g_{\mu\nu}(x)dx^2 - dxg_{\mu\nu}(x)) $$
    $$ dx = (g_{\mu\nu}(x)^2dx^2 - g_{\mu\nu}(x)dx g_{\mu\nu}(x))^{1 \over 2} $$
    と反重力と正規部分群の経路和となり、　
    $$ \pi(\chi, x) = i \pi (\chi, x)f(x) - f(x)  \pi (\chi, x) $$
    と基本群にまとまる。
    シュワルツシルト半径は、
    $$ ds^2 = -(1 - {r_s \over r})c^2 dt^2 + {1 \over {{1 - {r_s \over r}}}} 
    dr^2 + r^2 d \theta^2 + r^2 \sin^2 \theta d \psi^2 $$
    これは、カルーツァ・クライン空間と同型となる。
    一般相対性理論は、
    $$ R^{\mu\nu} + {1 \over 2}g_{ij}\Lambda = \kappa T^{\mu\nu} $$
    この多様体積分は、
    $$ \int \kappa T^{\mu\nu}\mathrm{dvol} = \int  \left(R^{\mu\nu} + 
    {1 \over 2}g_{ij}\Lambda\right)\mathrm{dvol}  $$
    ガンマ関数のおける大域的微分多様体は、これと同型により、
    $$ \int \Gamma \cdot {d \over {d \gamma}}{\Gamma dx_m} \le \int \Gamma
    dx_m + {d \over d \gamma}\Gamma $$
    $$ = \int \Gamma(\gamma)^{'}dx_m $$
    オイラーの定数の多様体積分は、
    $$ \int \left({{\int {1 \over x^s}}dx - \log x} \right)\mathrm{dvol} $$
    この解は、シュワルツシルト半径と同型より、
    $$ = e^f - e^{-f} \le {e^f + e^{-f}} $$
    次元の単位は多様体より、大域的微分多様体とオイラーの定数
    の多様体積分の加群分解は、オイラーの公式の三角関数の虚数の度解より、
    $$ {d \over df}F + \int C dx_m = 2 (\cos (i x \log x) 
    - i \sin (i x \log x)) $$
    すべては、大域的微分多様体の重力と反重力方程式に行き着き、
    $$ {d \over df}F \ge {d \over df}\int \int{1 \over (x \log x)^2}dx_m
    + {d \over df}\int \int{1 \over (y \log y)^{1 \over 2}}dy_m $$
    と求まる。


  ８種類の微分幾何での、それぞれの時間の固有値の動静は、
  ストークスの流体力学での、熱エネルギーの各流れの仕方で表されている。
  種数１の2種が種数０の球体に両辺で交わっていると、
一般相対性理論における重力理論は大域的微分多様体と積分多様体についての 単体量を求めるための空間の対数関数における不変性を記述するプランクスケール と異次元の宇宙におけるスケール、ウィークスケールと言われているanti-D-brane として、この２種類の計量をガンマ関数とベータ関数としてオイラーの定数 とそれによる連分数が微分幾何の量子化と数式の値を求めると同じという 予知と推測値から、広中平祐定理の４重帰納法のオイラーの公式による 多様体積分と同類値として、ペレルマン多様体がサーストン空間に成り立つ と同じく、広中平祐定理がヒルベルト空間で成り立つとしての、 これら２種の多様体の場の理論が、ゼータ関数がAdS5多様体で成り立つことと、 不変量として、ゼータ関数がこの3種の多様体の場の理論のバランスをとる
  理論として言えることである。
$$ ||ds^2|| = \lim_{x \to \infty} [\delta(x) \int \int \int \pi \left( \sum_{k=0}^{\infty}
		{{}^n\sqrt{p},x \over n} \right)
^{1 \over 2}d \tau]^{\mu\nu} $$
  This norm is component of fields in Hilbert manifold of space theorem rebuilt with
Mebius space in Gamma Function on Beta function fill into power of rout fundamental
group.

  　 Gravity of general relativity theory describe with Cutting of space
in being discatastrophed from Global differential and integral manifold
of scaled levetivity in plank scal and the other vector of universe
scale, these two scale inspectivity of Gamma and Beta function escourted
with these manifold experted in result of Differentail geometry of
quantum level manifold equal with Euler product and continue parameter,
moreover this product is relativity of Hironaka theorem in Four assembled
integral of Euler equation.

   \newcommand{\variint}{%
	\begin{picture}(15,30)
		\put(0,0){
			$ \iint $}
		\put(0,5){\line(15,0){15}}
	\end{picture} %
}

\newcommand{\variiint}{%
	\begin{picture}(15,15)
		\put(0,0){
			$ \iiint $}
		\put(0,5){\line(15,0){20}}
	\end{picture} %
}



 \newcommand{\alearletter}{%
	 \begin{picture}(10,20)
		 \put(0,0){
			 $ \square $}
		 \put(0,0){\line(1,1){10}}
	 \end{picture} %
	 }

 重力場理論の式は、
  Gravity equation is
  $$ \square = - {16 \pi G \over c^4}T^{\mu\nu} $$
  This equation quated with being logment of formula, and
  this formula divided with universe of number in prime zone,
  therefore, this dicided with varint equation is monotonicity of
  being composited with Weil's theorem united for Gamma function.
  この式は、対数を宇宙における数により求める素数分布論として、この
  大域的積分分断多様体がガンマ関数をヴェイユ予想を根幹とする単体量
  として決まることに起因する商代数として導かれる。
  $$ \scalebox{1}[2]{\alearletter} = 
  {{{{8 \pi G} \over {c^4}}T^{\mu\nu}}/ \log x} $$
  $$ {}^{t}\scalebox{1}[2]{\variiint}\mathrm{cohom D_{\chi}[I_m]}$$
  $$ = \oint {(px^n + qx + r)^{\nabla l}} $$
  $$ {d \over d l}L(x,y) = 2 \int ||\sin 2x||^{2}d \tau $$
  $$ {d \over d \gamma}\Gamma $$
  この関数は大域的微分多様体としてのアカシックレコードの合流地点として、
  タプルスペースを形成している。
  This function esterminate with acasic record of global differential
  manifolds.
  $$ = [i \pi (\chi, x), f(x)] $$
  それにより、この多様体は基本群をアカシックレコードの相対性としての
  存在論の実存主義から統合される多様体自身としてのタプルスペースの
  池になっている。
  And this manifold from fundemental group esterminate with
  also this manifold estimate relativity of acasic record.
  $$ {d \over d \gamma}\Gamma = \int C dx_m = \int(\int {1 \over x^s}dx
  - \log x)\mathrm{dvol} $$
  また、このアカシックレコードはオイラーの定数のラムダドライバー
  にもなっている。
  More also this record tupled with lake of Euler product.
  $$ \bigoplus({i \hbar^{\nabla}})^{\oplus L} $$
  $$ x^{{1 \over 2} + iy} = e^{x \log x} $$
  それゆえに、この定数はゼータ関数と微分幾何の量子化を因数にもつ
  素因子分解の式にもなっている。
  Therefore, this product also constructed with differential geometry
  of quantum level equation and zeta function.
  $$ = C $$ 
  そして、この関数はオイラーの定数から広中平祐定理による4重帰納法の
  オイラーの公式からの多様体積分へとのサーストン空間のスペクトラム関数
  ともなっている。
  And this function of Euler product respectrum of focus with
  Heisuke Hironaka manifold in four assembled of integral Euler equation.. 

  $$ C =  b_0 + \cfrac{c_1}{b_1 +
	  \cfrac{c_2}{b_2 +
	  \cfrac{c_3}{b_3 +
	  \cfrac{c_4}{b_4 + \cdots}}}} $$
  この方程式は指数による連分数としての役割も担っている。  
  This equation demanded with continued number of step function.
  $$ x^{{1 \over 2} + iy} = e^{x \log x} $$
  $$ (2.71828)^{2.828} = a^{a + b^{b + c^{c + d^{d \cdots }}}} $$
  $$ = \int e^f \cdot x^{1 - t}dx $$
  This represent is Gamma function in Euler product.
  Therefore this product is zeta function of global differential equation.

  $$ {d \over d f}F(v_{ij},h) = \int {e^{-f}\mathrm{dV}}
  [- \Delta v + \nabla_{i}\nabla_{j}
  v_{ij} - R_{ij}v_{ij} - v_{ij}\nabla_{i}\nabla_{j} + 2 <\nabla f,\nabla h>
  + (R + \nabla f^2)({v \over 2} - h)]$$
  $$ = {d \over d f}F = {2 {\int (R + \nabla_{i}\nabla_{j} f)^2} \over 
  { - (R + \Delta f)}}\mathrm{dm} $$
  これらの方程式は8種類の微分幾何の次元多様体として、そして、これらの
  多様体は曲平面による双対性をも生成している。 
  そして、このガウスの曲平面は、大域的微分多様体と微分幾何の量子化から
  素因数を形成してもいる。
  These equation are eight differential geometry of dimension calyement.
  And these calyement equation excluded into pair of dimension surface.
  This surface of Gauss function are global differential manifold,
  and differential geometry of quantum level.
  $$ F \ge {d \over d f}\int \int{1 \over (x \log x)^2}dx_m 
  + {d \over d f}\int \int{1 \over (y \log y)^{1 \over 2}}dy_m $$

  微分幾何の量子化はオイラーの定数とガンマ関数が指数による連分数
  としての不変性として素因数を形成していて、このガウスの曲平面
  による量子力学における重力場理論は、ダランベルシアンの切断多様体
  がこの大域的切断多様体を付加してもいる。
  Differential geometry of quantum level constructed with Euler product
  and Gamma function being discatastrophed from continued fraction style.
  Gravity equation lend with varint of monotonicity of level expresented
  from gravity of letter varient formula.
  これらの方程式は基本群と大域的微分多様体をエスコートしていもいて、
  ヴェイユ予想がこのダランベルシアンの切断方程式たちから
  輸送のポートにもなっている。ベータ関数とガンマ関数がこれらのフォームラ
  の方程式を放出してもいて、結果、これらの方程式は広中平祐定理の複素
  多様体とグリーシャ教授によるペレルマン多様体からサポートされてもいる。
  この2名の教授は、一つは抽象理論をもう一方は具象理論を説明としている。
  These equation escourted into Global differential manifold and
  fundemental equation.
  Weil's theorem is imported from this equation in gravity of letters.
  Beta function and Gamma function are excluded with these formulas.
  These equation comontius from Heisuke Hironaga of complex manifold
  and Gresha professor of Perelman manifold.
  These two professos are one of abstract theorem and the other of
  visual manifold theorem.

\section{Hilbert manifold in Mebius space \\
        this element of Zeta function on integrate of fields}

 Hilbert manifold equals with Von Neumann manifold, and this fields is concluded with Glassman manifold,
the manifold is dualty of twister into created of surface built.
These fields is used for relativity theorem and quantum physics.
$$ ||ds^2|| = \lim_{x \to \infty} [\delta(x) \int \int \int \pi \left( \sum_{k=0}^{\infty}
		{{}^n\sqrt{p},x \over n} \right)
^{1 \over 2}d \tau]^{\mu\nu} $$
  This norm is component of fields in Hilbert manifold of space theorem rebuilt with
Mebius space in Gamma Function on Beta function fill into power of rout fundamental
group.
And this result with $AdS_5$ space time in Quantum caos of Minkowsky of manifold in
abel manifold.
Gravity of metric in non-commutative equation is Global differential equation
conbuilt with Kaluza-Klein space. Therefore this mechanism is $T^{\mu\nu}$ tensor
is equal with $R^{\mu\nu}$ tensor. And This moreover inspect with laplace operator
in stimulate with sign of differential operator.
Minus of zone in Add position of manifold is Volume of laplace equation rebuilt with
Gamma function equal with summative of manifold in Global differential equation,
this result with setminus of zone of add summative of manifold, and construct with
locality theorem straight with fundamental group in world line of surface, this
power is boson and fermison of cone in hyper function.

$$ V(\tau) = [f(x),g(x)] \times [f^{-1}(x), h(x)] $$
$$ \Gamma(p,q) = \int e^{-x} x^{1 - t}dx $$
$$ = \beta(p,q) $$
$$ = \pi (f(\chi,x),x) $$
$$ ||ds^2|| = \mathcal{O}(x)[(f(x) \circ g(x))^{\mu\nu}]dx^{\mu}dx^{\nu} $$
$$ = \lim_{x \to \infty} \sum_{k=0}^{\infty}a_k f^k $$
$$ G^{\mu\nu} = {\partial \over \partial f}\int [f(x)^{\mu\nu} \circ G(x)^{\mu\nu}dx^{\mu}dx^{\nu}]^{\mu\nu}dm $$
$$ = g_{\mu\nu}(x)dx^{\mu}dx^{\nu} - f(x)^{\mu\nu}dx^{\mu}dx^{\nu} $$
$$ [i \pi(\chi,x),f(x)] = i \pi f(x) - f(x) \pi(\chi,x) $$
$$ T^{\mu\nu} = (\lim_{x \to \infty} \sum_{k=0}^{\infty}\int \int [V(\tau) \circ S^{\mu\nu}(\chi,x)]dm)
^{\mu\nu}dx^{\mu}dx^{\nu} $$
$$ G^{\mu\nu} = R^{\mu\nu}T^{\mu\nu} $$
$$ \begin{vmatrix} D^m & dx \\
dx & \sigma^m \end{vmatrix} \begin{vmatrix} \cos \theta & - \sin \theta \\
\sin \theta & \cos \theta \end{vmatrix} \begin{vmatrix} x \\
y \end{vmatrix} = \begin{pmatrix} 1 & 0 \\ 
0 & - 1 \end{pmatrix}^{1 \over 2} $$
$$ \sigma^m \begin{bmatrix} \delta(x) & - 1 \\
1 & \epsilon(x) \end{bmatrix}^{1 \over 2} = \begin{pmatrix} i & 0 \\
0 & - i \end{pmatrix} $$
$$ V(M) = {\partial \over \partial f} ({}^N \int [f \smallsetminus M]^{\oplus N})^{\mu\nu} dx^{\mu}dx^{\nu} $$
$$ V(M) = \pi(2 \int \sin^2 dx)\oplus {d \over df}F^M dx_m $$
$$ \lim_{x \to \infty}\sum_{k=0}^{\infty} a_k f^k = \int (F(V)dx_m)^{\mu\nu}dx^{\mu}dx^{\nu} $$
$$ \bigoplus_{k=0}^{\infty}[ f \smallsetminus g] = \vee(M \wedge N) $$
$$ \pi_1 (M) = e^{-f 2 \int \sin^2 x dm} + O(N^{-1}) $$
$$ = [i \pi(\chi,x),f(x)] $$
$$ M \circ f(x) = e^{-f \int \sin x \cos x dx_m} + \log (O(N^{-1})) $$

 Non-Symmetry space time.
 $$ {d \over dL}V(\tau) = {d \over df}\int \int_M {1 \over (x \log x)^2}dx_m + {d \over df}\int \int_M
 {1 \over (y \log y)^{1 \over 2}}dy_m $$
 $$ \epsilon S(\nu) = \square_v \cdot {\partial \over \partial \chi}({}^5 \sqrt{\wedge g^2})d\chi $$
 Differential Volume in $AdS_5$ graviton of fundamental rout of group.
 $$ \wedge (F^m_t)^{''} = {1 \over 12}g_{ij}^2$$
 Quarks of other dimension.
 $$ \pi(V_{\tau}) = {e^{-(\sqrt{\pi \over 16}\log x)}}^{\delta} \times {1 \over (x \log x)} $$
 Universe of rout, Volume in expanding space time.
 $$ {d \over dt}(g_{ij})^2 = {1 \over 24}(F^m_t)^2 $$
 $$ m^2 = 2 \pi T \left({26 - D_n \over 24}\right) $$
 This quarks of mass in relativity theorem, and fourth of universe in three manifold of one dimension surface,
 and also this integrate of dimension in conbuilt of quarks.
 $$ g_{ij} \wedge \pi(\nu_{\tau}) = e^{-2 \pi T |\psi|}[\eta_{\mu\nu} + \bar{h}_{\mu\nu}(x)]dx^{\mu}dx^{\nu} 
 + T^2 d \psi^2 $$
 Out of rout in $AdS_5$ space time.
 $$ ||ds^2|| = g_{ij}\wedge \pi(\nu_v) $$
 $AdS_5$ norm is fourth of universe of power in three manifold out of rout.


  Sphere orbital cube is on the right of hartshorne conjecture by the right equation,
this integrate with fourth of piece on universe series, and this estimate with one of three manifold.
Also this is blackhole on category of symmetry of particles.
These equation conbuilt with planck volume and out of universe is phologram field, this field construct with
electric field and magnitic field stand with weak electric theorem.
This theorem is equals with abel manifold and $AdS_5$ space time.
Moreover this field is anti-brane and brane emerge with gravity and antigravity equation restructed with.
Zeta function is existed with Re pole of ${1 \over 2}$ constance reluctances.
This pole of constance is also existed with singularity of complex fields.
Von Neumann manifold of field is also this means.
$$ ||ds^2|| = \lim_{x \to \infty} [\delta(x)\int \int \int \pi \left( \sum_{k=0}^{\infty}{{}^n\sqrt{p},x \over n}\right)
^{1 \over 2}d \tau]^{\mu\nu} $$
$$ e^{i \theta} = \cos \theta + i \sin \theta $$
$$ e^{x \log x} = x^{{1 \over 2} + iy}, x \log x = \log (\cos \theta + i \sin \theta) $$
$$ = \log \cos \theta + i \log \sin \theta $$
This equation is developed with frobenius theorem activate with logment equation.
This theorem used to
$$ \log(\sin \theta + i \cos \theta)= \log(\sin \theta - i \cos \theta) $$
$$ \log \left({\sin \theta \over i \cos \theta}\right) = -2 R_{ij}, {d \over dt}g_{ij}(t) =  - 2 R_{ij} $$
$$ \mathcal{O}(x) = {\zeta(s) \over \sum\limits_{k=0}^{\infty}a_k f^k} $$
$$ \mathrm{Im} f = \mathrm{ker} f, \chi(x) = {\mathrm{ker}f \over \mathrm{Im}f} $$ $$ H(3) = 2, \nabla H(x) = 2, 
 \pi(x) = 0 $$
 This equation is world line, and this equation is non-integrate with relativity theorem of rout constance.

$$ [f(x)] = \infty, ||ds^2|| = \mathcal{O}(x)[\eta_{\mu\nu} + \bar{h}_{\mu\nu}(x)]dx^{\mu}dx^{\nu} + T^2 d^2 \psi $$
$$ T^2 d^2 \psi = [f(x)] , T^2 d^2 \psi = \lim_{x \to 1}\sum_{k=0}^{\infty}a_k f^k $$

 These equation is equals with M theorem. For this formula with zeta function into one of universe in four dimensions.
Sphere orbital cube is on the surface into $AdS_5$ space time, and this space time is Re pole and Im pole of 
constance reluctances.


  Gauss function is equals with Abel manifold and Seifert manifold. Moreover this function is infinite time,
and this energy is cover with finite of Abel manifold and Other dimension of seifert manifold on the surface.


  Dilaton and fifth dimension of seifert manifold emerge with quarks and Maxwell equation, this power
is from dimension to flow into energy of boson. 
This boson equals with Abel manifold and $AdS_5$ space time, and this space is created with Gauss function.


  Relativity theorem is composited with infinite on D-brane and finite on anti-D-brane, This space time
restructed with element of zeta function.
Between finite and infinite of dimension belong to space time system, estrald of space element have with 
infinite oneselves. Anti-D-brane have infinite themselves, and this comontend with fifth dimension of AdS5 
have for seifert manifold, this asperal manifold is non-move element of anti-D-brane and move element 
of D-brane satisfid with fifth dimension of seifert algebra liner.
Gauss function is own of this element in infinite element.
And also this function is abel manifold.
Infinite is coverd with finite dimension in fifth dimension of AdS5 space time.
Relativity theorem is this system of circustance nature equation.
AdS5 space time is out of time system and this system is belong to infinite mercy.
This hiercyent is endrol of memolite with geniue of element.
Genie have live of telomea endore in gravity accesorlity result.
AdS5 space time out of over this element begin with infinite assentance.
Every element acknowlege is imaginary equation before mass and spiritual envy.

  $$ T^2 d^2 \psi = \lim_{x \to 1}\sum_{k=0}^{\infty}a_k f^k $$
  $$ \lim_{x \to 1}\sum_{k=0}^{\infty}a_k f^k = [T^2 d^2 \psi] $$

  $$ {d \over dL}V(\tau)= {d \over df}\int \int_M({}^5\sqrt{x^2})d\Lambda 
  + {d \over df}\int \int_M {}^N({}^3\sqrt{x})^{\oplus N}d\Lambda $$
  $$ {}^M(\vee(\wedge f \circ g)^N)^{1 \over 2} = {d \over df}\int \int_M{1 \over (x \log x)^2}dx_m
  + {d \over df}\int \int_M{1 \over (y \log y)^{1 \over 2}}dy_m $$
  $$ ||ds^2|| = \mathcal{O}(x)[\eta_{\mu\nu} + \bar{h}_{\mu\nu}(x)]dx^{\mu}dx^{\nu}
  + T^2 d^2 \psi $$
  $$ \mathcal{O}(x) = e^{-2\pi T|\psi|} $$
  $$ G^{\mu\nu} = R_{\mu\nu} T^{\mu\nu} $$
  $$ = - {1 \over 2}\Lambda g_{ij}(x) + T^{\mu\nu} $$

  Fifth dimension of eight differential structure is integrate with one geometry element,
Four of universe also integrate to one universe, and this universe represented with symmetry formula.
Fifth universe also is represented with seifert manifold peices.
All after universe is constructed with six element of circuatance.
This aquire maniculate with quarks of being esperaled belong to.

\section{Imaginary equation in AdS5 space time create with \\
       dimension of symmetry}

  D-brane and anti-D-brane is composited with all of series universe emerged for one geometry of 
dimension, this gravity of power from D-brane and anti-D-brane emelite with ancestor.
Seifert manifold is on the ground of blackhole in whitehole of power pond of senseivility.
Six of element of quarks and universe of pieces is supersymmetry of mechanism resolved with

hyper symmetry of quarks constructed to emerge with darkmatter.
This darkmatter emerged with big-ban of heircyent in circumstance of phenomena.

  D-brane and anti-D-brane equations is
  $$ {d \over dL}V(\tau) = {d \over df}\int \int_M {1 \over (x \log x)^2}dx_m
  + {d \over df}\int \int_M {1 \over (y \log y)^{1 \over 2}}dy_m $$
  $$ C = \int \int {1 \over (x \log x)^2}dx_m $$

  Euler constance is quantum group theorem rebuild with projective space involved with.


   虚数方程式は、反重力に起因するフーリエ級数の励起を生成する。それは、人工知能を生み出す、
   ５次元時空にも、この虚数方程式は使われる。AdS5の次元空間は、反ド・シッター時空の
   D-braneとanti-D-braneのcomformal場を生み出す。ホログラフィー時空は、この量子起因によるものである。
   2次元曲面によるブラックホールは、ガンマ線バーストによる５次元時空の構造から観測される。
   空間の最小単位によるプランク定数は、宇宙の大域的微分多様体から導き出される、AdS5の次元空間の
   準同型写像を形成している。これは、最小単位から宇宙の大きさを導いている。最大最小の方程式は、
   相加相乗平均を形成している。時間と空間は、宇宙が生成したときから、宇宙の始まりと終わりを既に生み出している。
   宇宙と異次元から、ブラックホールとホワイトホールの力がわかり、反重力を見つけられる。オイラーの定数は、
   この量子定数からわかる。虚数の仕組みはこの量子スピンの産物である。
   オイラーの定数は、 この虚時空の斥力の現存である。それは、非対称性理論から導かれる。
   不確定性原理は、AdS5のブラックホールとホワイトホールを閉３次元多様体に統合する５次元時空から求められる。
   位置と運動エネルギーが、空間の最小単位であるプランク定数を宇宙全体にする微細構造定数からわかり、
   面積確定から、アーベル多様体を母関数に極限値として、ゼータ関数をこの母関数に不変式として、
   ２種類ずつにまとめる４種類の宇宙を形成する８種類のサーストンの幾何化予想から導き出される。
   この閉３次元多様体は、ミラー対称性を軸として、６種類の次元空間を一種類の異次元宇宙と同質ともしている。
   複素多様体による特異点解消理論は、この原理から求められる。この特異点解消理論は、２次元曲面を３次元多様体
   に展開していく、時空から生成される重力の密度を反重力と等しくしていく時間空間の４次元多様体と虚時空
   から求められる。
   ヒルベルト空間は、フォン・ノイマン多様体とグラスマン多様体をこのサーストンの幾何化予想を場の理論
   既定値として形成される。
   この空間は、ミンコフスキー時空とアーベル多様体全体を表している。そして、この空間は、球対称性を
   複素多様体を起点として、大域的トポロジーから、偏微分を作用素微分として時空間をカオスから
   ずらすと５次元多様体として成立している。
   これらより、３次元多様体に２次元射影空間が異次元空間として、AdS5空間を形成される。
   偏微分、全微分、線形微分、常微分、多重微分、部分積分、置換微分、大域的微分、単体分割、
   双対分割、同調、ホモロジー単体、コホモロジー単体、群論、基本群、複体、マイヤー・ビーートリス完全系列、
   ファン・カンペンの定理、層の理論、コホモルティズム単体、CW複体、ハウスドロフ空間、線形空間、位相空間、
   微分幾何構造、モースの定理、カタストロフの空間、ゼータ関数系列、球対称性理論、スピン幾何、
   ツイスター理論、双対被覆、多重連結空間、プランク定数、フォン・ノイマン多様体、グラスマン多様体、
   ヒルベルト空間、一般相対性理論、反ド・シッター時空、ラムダ項、D-brane,anti-D-brane,コンフォーマル場、
   ホログラム空間、ストリング理論、収率による商代数、ニュートンポテンシャルエネルギー、剛体力学、
   統計力学、熱力学、量子スピン、半導体、超伝導、ホイーストン・ブリッチ回路、非可換確率論、Connes理論、
   これら、演算子代数を形成している、微分・積分作用素が、ヒルベルト空間に存在している多様体の特質を
   全面に押し出して、いろいろな多様体と関数そして、群論を形作っている。


$$ \begin{vmatrix} D^m & dx \\
dx & \sigma^m \end{vmatrix} \begin{vmatrix} \cos \theta & - \sin \theta \\
\sin \theta & \cos \theta \end{vmatrix} \begin{vmatrix} x \\
y \end{vmatrix} = \begin{pmatrix} 1 & 0 \\ 
0 & - 1 \end{pmatrix}^{1 \over 2} $$
$$ \sigma^m \begin{bmatrix} \delta(x) & - 1 \\
1 & \epsilon(x) \end{bmatrix}^{1 \over 2} = \begin{pmatrix} i & 0 \\
0 & - i \end{pmatrix} $$
$$ (D^m,dx)\cdot (\cos \theta, \sin \theta) \times (dx, \partial^m)\cdot (\cos \theta, \sin \theta) $$

$$ \begin{pmatrix} 1 && 0 \\
   0 && i \end{pmatrix} \beta(x,\theta) \begin{pmatrix} x \\
   y \end{pmatrix} $$
   $$ = \begin{pmatrix} i && 0 \\
   0 && - 1 \end{pmatrix} $$
   $$ l = {\sqrt{\hbar G \over c^3}}, \sigma^m \cdot \begin{pmatrix} \delta(x) && - 1 \\
   1 && \epsilon (x) \end{pmatrix}^{1 \over 2} = \begin{pmatrix} i && 0 \\
   0 && - i \end{pmatrix} $$
   String theorem also imaginary equation of pole with the fields of manifold.
   $$ \left({\partial \over \partial \tau}f(x,y,z)\right)^{3'} = A^{\mu\nu} $$
   $$ {d \over dt}g_{ij}(t) = - 2 R_{ij} $$
   Rich equation is all of top formula in integrate theorem.
      $$ = (D^m,dx) \cdot (\cos \theta, \sin \theta) \times (dx, \partial^m)\cdot (\cos \theta, \sin \theta) $$
      Dalanverle equation equals with abel manifold.






\title{慣性の法則と、回転体による反重力 \\
タイミングとともに、別ベクトルとして重力を加える法則}
\author{Masaaki Yamaguchi}


   慣性の法則の慣性力が働くときと同時に、
   前もって、タイミングを計り、慣性力とともに、
   逆向きのベクトルとして、慣性力の力と同じエネルギーの
   重力を生成して、慣性力の影響を無くす。この力を反重力という。
   慣性力と逆向きのベクトルとして、反重力として、
   慣性力に加えて、別のベクトルとする。

   反重力の本質は、重力のベクトルの違いであり、
   核エネルギーが特殊相対性理論と一般相対性理論を使って、
   放射性物質吸収体に、この原理を使い、核エネルギーをリサイクルしている。
   これを宇宙人は知っていて、UFOにこの原理を使っている。

   Higgs場は、重力であり、周りを反重力が覆っている。
   一般相対性理論のベクトルの向きを変えた力の装置が、
   ラムダ・ドライバであり、反重力である。

   重力は熱エネルギーである。ベクトルの向きが違う。
   反重力は吸収体である。
   慣性の法則が働いているときに、これが手に入ったが、離れたになり、
   重力をベクトルの向きを変えて、同じエネルギーを加えた場合、
   量子コンピュータが、そのときの差分を使って、反重力が慣性の法則
   を変えて、宙に浮くUFOにしている。 

   まとめると、慣性の法則のときに、重力のベクトルを違う向きとして、
   加えると、手に入ったが、離れたになり、これが反重力である。 

   UFOが現れる前に、ヘリコプターが現れて、そのあとに、UFOが実際に現れる。
   慣性の法則のしばらくして、反重力として、UFOとして、
   宇宙人が人に教えている。

   反重力の生成は、回転体の電磁気力からくるちからであり、ローレンツ力と
   同じでもあり、機体の内部で生成される力であり、そのために、放射性
   物質吸収体を結晶石として、このエネルギーからくる放射力を吸収する
   ために、結晶石に反重力の放射線を吸収する。その上に、外部の
   慣性力を内部の反重力で、手に入ったら、離れたをしている。


\title{反重力発生装置の各構成物質}
\author{Masaaki Yamaguchi}


  原子振動子　セシウム Cs


  形状記憶合金 (Fe・Co60・Pt・Al)$H_2 SO_4, Al_2 O_3$


  反重力発生器 Pr(パラジウム合金)電磁場生成He,$H_2$Mg,Al(摩擦熱)


  結晶石 Co60,Pt,Prの反陽子


  緩衝剤 慣性力感知器 有機化合物とS化合物(シクロアルカン$C_n H_{2n}$
  方位と位置探知機にも使う


  量子コンピュータの量子素子 Pt,Ag,Au（電子と電流のエネルギー経路)半導体に
  使う
  回路の合成演算子 


  $$ x^y = {1 \over {y^x}} \to $$

  $$ \pi(\chi,x) = [i\pi(\chi,x),f(x)] $$
  $$ {d \over df}F = F^{{f}^{'}} $$


  ディスプレイの電磁迷彩 (Al,Mg) $\to$ S合成子 $Si O_2 $
  プリズム C

\end{document}
